% 本文件由 LaTeX 排版 Agent 根据有效 translation.json 自动生成。
% author=John Chrysostom; work_id=2310; title=迦拉达书讲道集

\chapter{迦拉达书讲道集}

% structure: p0001 source=h1 target=omitted reason=page-title-already-rendered-by-parent

第一章\newline{}第二章\newline{}第三章\newline{}第四章\newline{}第五章\newline{}第六章\par

\SourceRecord{Translated by Gross Alexander.；Nicene and Post-Nicene Fathers, First Series；Vol. 13.；Edited by Philip Schaff.；Buffalo, NY: Christian Literature Publishing Co.,；1889.；<http://www.newadvent.org/fathers/2310.htm>.}

% structure: p0001 source=h1 target=section reason=page-title

\section{迦拉达书讲道第一讲}

第一至三节 \BibleRef{迦 1:1}\par

\BibleQuote{保禄宗徒，（不是由于人，也不是借着人，而是借着耶稣基督和那使他从死者中复活的天主圣父，）以及同我在一起的众弟兄，致书给迦拉达各教会：愿恩宠与平安由天主圣父和我们的主耶稣基督赐与你们。}\par

开场白充满激昂和高昂的精神，不仅开场白，而且可以这样说，整封书信都是如此。因为总是以温和的态度对待门徒，即使他们需要严厉的时候，这不是老师的行为，而是败坏者和敌人的行为。因此，我们的主虽然通常温和地对待门徒，但也不时使用严厉的语言，一会儿祝福，一会儿责备。先前对伯多禄说了：\BibleQuote{西满·巴尔约拿，你是有福的，}\BibleRef{玛 16:17}并应许将教会的根基建立在祂的宣认上，不久后祂说：\BibleQuote{退到我后面去，撒殚！你是我的绊脚石。}\BibleRef{玛 16:23}再后来，祂又说：\BibleQuote{你们还是不明白吗？}\BibleRef{玛 15:16}祂使他们感到敬畏，这从若望的话可以看出：当他们看见祂与撒玛黎雅妇人谈话时，虽然提醒祂进食，却没有一个人敢说：\BibleQuote{你找什么？你为什么和她谈话？}\BibleRef{若 4:27}这样受教，效法他老师的脚步，保禄根据门徒的需要调整言辞，时而使用刀和烧烙，时而施以温和的疗法。对格林多人他说：\BibleQuote{你们愿意怎样？要我带着棍棒到你们那里去，还是要我怀着慈爱和温柔的心？}\BibleRef{格前 6:21}但对迦拉达人，他说：\BibleQuote{无知的迦拉达人！}\BibleRef{迦 3:1}他不止一次，而是第二次使用了这种责备，在结束时他带着责备暗示说：\BibleQuote{从今以后，不要有人再烦扰我。}\BibleRef{迦 6:17}但他又用言语安慰他们：\BibleQuote{我的孩子们！我再次为你们受产痛。}\BibleRef{迦 4:19}其他许多例子也是如此。\par

这封书信充满愤怒的情绪，这一点甚至初次阅读时每个人都显而易见；但我必须解释他为何对门徒发怒。原因不可能轻微或不重要，否则他不会如此激烈。因为因小事而恼怒是小气、乖戾、暴躁之人的特征，就像在重大事情上灰心丧气是更为迟钝和懒惰之人的特征一样。保禄不是这样的人。那么，是什么冒犯激怒了他？这是一件严重而重大的事，使他们全都与基督疏远，正如他随后所说：\BibleQuote{我保禄告诉你们：如果你们接受割损，基督对你们就没有益处。}\BibleRef{迦 5:2}又说：\BibleQuote{你们这些要靠法律成义的人，是与恩宠隔绝了。}\BibleRef{迦 5:4}那么，这是什么？必须更清楚地解释。一些信了主的犹太人，受犹太主义先入之见的束缚，同时又陶醉于虚荣，渴望为自己获得教师的尊严，来到迦拉达人那里，教导他们遵守割损、安息日和月朔是必要的，而保禄废除这些事是不能容忍的。因为他们说：伯多禄、雅各伯和若望，这些宗徒之长和基督的同伴，并没有禁止这些事。事实上，他们并没有禁止这些事，但这并不是在发表确定的教义，而是迁就犹太信徒的软弱。保禄在向外邦人传教时不需要这种迁就，但他在犹太地时，自己也采用了这种做法。但这些欺骗者，隐瞒了保禄及其弟兄迁就的原因，误导了单纯的人，说他不应该被容忍，因为他昨天才出现，而伯多禄和他的同伴从一开始就在；说他是宗徒的门徒，而他们却是基督的门徒；说他只是一个人，而他们却是许多人，是教会的柱石。他们还指责他虚伪，说这位禁止割损的人，在别处却遵行这仪式，向你们宣讲一套，向别人宣讲另一套。\par

保禄既看见整个迦拉达人民处于兴奋状态，火焰对着他们的教会燃起，大厦摇摇欲坠，充满了义怒和沮丧的混合情感，（他以\BibleQuote{我恨不得现今在你们当中，改变我的声音}\BibleRef{迦 4:20}表达了这一点，）便写下这封书信作为对这些指控的回应。从一开始他的目的就在于此，因为那些诋毁他声誉的人说，其他人是基督的门徒，而这个人却是\Quote{宗徒}的门徒。因此他这样开头：\BibleQuote{保禄，宗徒，不是从人来的，也不是借着人来的}。因为，这些骗子（正如我之前所说）曾说过，这个人是所有宗徒中最末的一位，并且是由他们教导的，因为伯多禄、雅各伯和若望既是最先蒙召的，又在门徒中占有首位，并且还亲自由基督本人领受了他们的教义；因此，服从他们更胜过服从这个人；并且他们不禁止割礼也不禁止遵守法律。通过这种和类似的言论，以及贬低保禄、抬高其他宗徒的荣誉（尽管并不是为了赞扬他们，而是为了欺骗迦拉达人），他们诱导迦拉达人不合时宜地恪守法律。因此他的开头是恰当的。当他们诋毁他的教义，说它来自人，而伯多禄的教义来自基督时，他立即直接处理这一点，宣称自己是\BibleQuote{不从人来的，也不借着人}的宗徒。是阿纳尼雅为他施洗，但不是阿纳尼雅使他脱离错路并引入信仰；而是基督本人从天上发出那奇妙的声音，用这声音将他网住。至于伯多禄和他的兄弟、若望和他的兄弟，基督是在海边行走时召唤了他们\BibleRef{玛 4:18}，但保禄是在他升天之后\BibleRef{宗 9:3}。正如那些人不需要第二次召唤，立刻撇下网和一切所有跟从他一样，这个人也在第一次蒙召时奋勉前进，受洗后就立即与犹太人展开了无情的战争。在这方面，他主要胜过其他宗徒，如他所说：\BibleQuote{我比他们众人更劳苦}\BibleRef{格前 15:10}；然而此刻，他没有提出这样的要求，而是满足于与他们平起平坐。实际上，他的主要目的不是为自己建立任何优越性，而是推翻他们错误的根基。\Quote{不从人来的}这一点适用于所有的宗徒，因为福音的根源和起源是天主的；但\Quote{不借着人}则是宗徒们特有的；因为天主召唤他们，不是借由人的代理，而是借由他自己。\par

但是为什么他不谈论自己的蒙召，反而谈论他的宗徒职分，并说\Quote{保禄}蒙召\Quote{不是由人}？因为全部问题就在于此；因为他们说教师职务是由人，即由宗徒们，委托给他的，因此他应该服从他们。但这不是由人委托给他的，路加在以下话中宣告：\BibleQuote{他们事奉主并禁食的时候，圣神说：\InnerQuote{你们为我分开巴尔纳伯和扫禄。}}\BibleRef{宗 13:2}\par

从这段经文可以明显看出，圣子和圣神的能力是同一的，因为被圣神派遣，他就说被基督派遣。这在另一处也表现出来，从他将在米肋托向长老们说的话中归给圣神的事，可以看到：\BibleQuote{你们要谨慎自己和全群，圣神立你们为监督}\BibleRef{宗 20:28} 但在另一封书信中他说：\BibleQuote{天主在教会中设立的，第一是宗徒，第二是先知，第三是教师}\BibleRef{格前 12:28} 因此，他毫不区分地将属于圣神的事归给天主，将属于天主的事归给圣神。在此，他也用\Quote{借着耶稣基督和天主父}这句话堵住了异端者的口；因为，他们说\Quote{借着}这个词用于圣子时暗示了次级性，看看他做了什么：他把它用于圣父，从而教导我们不要给不可言说的本性制定法则，也不要定义属于圣父和圣子的神性等级。因为在\Quote{借着耶稣基督}这句话之后，他补充了\Quote{和天主父}；因为若只提及圣父时他使用了\Quote{借着祂}这个短语，他们可能会诡辩说它特别适用于圣父，因为圣子的行为要归给祂。但他同时提及圣子和圣父，并使他的语言适用于两者，没有留下任何吹毛求疵的余地。他这样做，并不是将圣子的行为归给圣父，而是表明这个表达方式并不暗示本质的区别。此外，那些从圣洗圣事格式——我们受洗归于父、子和圣神的名下——得出次级性观念的人，现在还能说什么呢？因为如果圣子因被列在圣父之后而成为次级者，那么看看在我们面前的这段经文中，宗徒从基督开始然后提到圣父，他们又会说什么呢？——但我们甚至不要说出这样的亵渎之言，在与他们的争论中不要偏离真理；相反，即使他们狂怒千万次，我们也要保持应有的敬畏尺度。既然因基督先被提到而论证圣子大于圣父是极大的疯狂和邪恶，那么我们也不敢因圣子在圣洗格式中被放在圣父之后而认为圣子次于圣父。\par

\BibleQuote{那使他从死者中复活的}\par

啊，\Person{保禄}，你本意是要引导这些犹太派归向信德，却为什么不提你致\BibleBook{斐理伯}{人书}中所论的任何那些伟大而光辉的论点，例如\BibleQuote{祂虽具有天主的形体，却没有以自己与天主同等为应紧持不舍的}\BibleRef{斐 2:6}，或是你后来在\BibleBook{希伯来}{书}中所宣布的\BibleQuote{祂是光荣的辉耀，是祂本体的真像}\BibleRef{希 1:3}，或是雷霆之子在他的福音开端所发出的声音：\BibleQuote{在起初已有圣言，圣言与天主同在，圣言就是天主}\BibleRef{若 1:1}，或是\Person{耶稣}自己屡次向犹太人宣布的\BibleQuote{祂的能力和权柄与圣父相等}\BibleRef{若 5:19}？你却忽略这一切，只提到祂降生成人的安排，提出祂的十字架和死亡？\Person{保禄}会回答：\Quote{是的。}因为如果这篇讲论是针对那些对基督持有卑下观念的人，那就应该提到那些事；但既然扰乱来自那些害怕放弃法律会受惩罚的人，他就提到那排除法律一切需要的事，即通过十字架和复活所赐予众人的恩惠。如果说\Quote{在起初已有圣言}，以及\Quote{祂具有天主的形体，使自己与天主平等}之类，固然显出圣言的天主性，但对当前的问题却毫无助益。而加上\Quote{谁使祂从死者中复活}却极为贴切，因为这提醒我们最大的恩惠，而人通常对谈论天主的威严，不如对阐述赐予人类的恩惠更感兴趣。因此，他舍弃前者，论述赐予我们的恩惠。\par

但这里异端者侮辱地喊道：\Quote{看，圣父使圣子复活！}因为他们一旦受了感染，就故意对所有更崇高的教义充耳不闻；他们只抓住并坚持那些较为卑微的、用较低词语表达的内容，无论是由于圣子的人性，或是为尊崇圣父，或是为其他暂时的目的，他们凌辱——我不说圣经——而是他们自己。我极想问这样的人：你们为什么这样说？你们希望证明圣子软弱无力，连\Emph{一个}身体也不能复活吗？不，实在，对祂的信德使那些信祂的人的影子也能叫死人复活。\BibleRef{宗 5:15} 那么，信祂的人，虽是可死的，却凭他们属世身体的影子，以及触摸过这身体的衣服，就能复活死人，而祂自己却不能复活自己吗？这不是明显的疯狂、极大的愚蠢吗？你们没有听过祂的话：\BibleQuote{拆毁这座圣殿，三天之内我要把它重建起来}\BibleRef{若 2:19} 又说：\BibleQuote{我有权舍掉我的生命，也有权取回它来}\BibleRef{若 10:18} 那么，为什么圣父被说成使祂复活，也做了圣子自己做的其他事？这是为尊崇圣父，并出于对听者软弱的怜悯。\par

\BibleQuote{还有与我同在的众弟兄。}\par

为什么他在其他时候写信从未加上这句话？因为他只放自己名字，或两三个人的名字，但这里提到了全体人数，而没有指名道姓。\par

那么他这样做是什么原因？\par

他们诽谤说他的讲论是独特的，想在基督信仰的教导中引入新奇。因此，为了消除他们的猜疑，并表明有许多人支持他的教义，他把自己与\Quote{弟兄们}联合起来，以表明他所写的是得到他们同意的。\par

\BibleQuote{致迦拉达的各教会。}\par

由此可见，谬误的火焰不仅蔓延了一两个城市，而是整个\Place{迦拉达}民族。也请思考，在\Quote{致迦拉达的各教会}这一短语中所包含的强烈愤慨：他没有说\Quote{致蒙爱的}或\Quote{致蒙圣化的}，并且省略了所有亲爱或尊敬的称呼，只是把他们当作一个团体来称呼，没有加上\Quote{天主的众教会}，而只是说\Quote{迦拉达的众教会}，这强烈表达了深切的忧虑和悲伤。在这里一开始，如同其他地方一样，他攻击他们的不当行为，因此给他们\Quote{教会}的名称，是为了使他们羞愧，并使他们归于合一。因为分裂成许多派别的人不能恰当地拥有这个称号，因为\Quote{教会}的名称是和谐与一致的名字。\par

\BibleQuote{愿恩宠与平安由天主我们的父和我们的主耶稣基督赐予你们。}\par

他总认为这是不可或缺的，特别是在这封致\Place{迦拉达}人书中；因为他们正处于从恩宠坠落的危险中，他祈祷他们能重新获得恩宠；又因为他们已与天主为敌，他恳求天主使他们恢复同样的平安。\par

\Quote{天主父}\par

这里再次直接驳斥了异端者，他们声称若望在福音的开端，当时他说\Quote{圣言就是天主}时，使用了无冠词的\OriginalTerm{\GreekText{Θεὸς}}{\GreekText{Θεὸς}}一词，以暗示圣子天主性的低劣；而保禄在说圣子\Quote{具有天主的形体}时，也并非指圣父，因为\OriginalTerm{\GreekText{Θεὸς}}{\GreekText{Θεὸς}}一词无冠词。因为在此处保禄说\OriginalTerm{从天主父}{\GreekText{ἀπὸ} \GreekText{Θεοῦ} \GreekText{Πατρός}}，而不是\OriginalTerm{从天主}{\GreekText{ἀπὸ} \GreekText{τοῦ} \GreekText{Θεοῦ}}，他们又能作何辩解呢？他称天主为\Quote{父}，并非出于对他们的慈爱，而是严厉的斥责，并提醒他们成为儿子的根源，因为这份尊荣并非通过法律赐予他们，而是通过重生的洗涤。因此，即使在开篇，他也处处散布天主仁慈的痕迹，我们可以想象他这样说：\Quote{你们这些不久前还是奴隶、仇敌和外人的人，怎么突然有权称天主为你们的父呢？赐予你们这种关系的不是法律；那么你为什么离弃那使你们如此亲近天主的那一位，而回到你们的监护人那里呢？}\par

但是，圣子的名号，以及圣父的名号，本已足以向他们宣告这些祝福。如果我们留心思考主耶稣基督的名号，这一点就会显明；因为经上说：\BibleQuote{你要给祂起名叫耶稣，因为祂要把自己的民族从他们的罪恶中拯救出来。}\BibleRef{玛 1:21}而\Quote{基督}的称号则使人想起圣神的傅油。\par

第4节。\BibleQuote{祂为我们的罪恶舍弃了自己。}\BibleRef{迦 1:4}\par

由此可见，祂所承担的职份是自由且非被迫的；祂是由自己而非由他人交付的。因此，不要让若望的话——\BibleQuote{天主竟这样爱了世界，甚至赐下了自己的独生子}\BibleRef{若 3:16}——引导你们贬低独生子的尊位，或由此推断祂仅仅是人。因为圣父被称为赐下了祂，并非暗示圣子的职份是奴仆式的，而是教导我们这在圣父眼中是好的，正如保禄在紧接的上下文中也表明的：\Quote{按照我们的天主和父的旨意。} 祂不是说\Quote{藉着命令}，而是\Quote{按照旨意}，因为圣父与圣子旨意合一，圣子所愿意的，圣父也愿意。\par

\Quote{为了我们的罪恶，}宗徒说道；我们曾以万般邪恶刺伤自己，并应受最严厉的惩罚；法律非但没有解救我们，反而定我们的罪，使罪更显明，既无力使我们脱离罪恶，也无法平息天主的义怒。但是天主子使这不可能成为可能，因为祂赦免了我们的罪恶，祂使我们从敌对转为朋友，祂白白赐给我们无数其他祝福。\par

第4节。\BibleQuote{为要把我们从这邪恶的现今世代中救出来。}\BibleRef{迦 1:4}\par

另一类异端者抓住保禄的这些话，曲解他的见证来指控现世生命。瞧，他们说，他称这个现世为邪恶的，请问\Quote{世界} [age] \OriginalTerm{世界（时代）}{\GreekText{αἴων}}除了以日时计算的时间，还意味着什么？难道日子的区分和太阳的运行是邪恶的吗？即使一个人被带至完全无理的地步，也不会这样断言。他们说，\Quote{但是}他们说，\Quote{他称之为邪恶的，不是\InnerQuote{时间}，而是现世\InnerQuote{生命}。}事实上，这些话本身并非如此说；但异端者不满足于字面，从中捏造指控，却为自己提出一种新的解释模式。因此，至少他们必须让我们提出我们的解释，更何况这解释既虔诚又合理。那么，我们断言，邪恶不能是善的原因，然而现世生命却产生上千种奖品和赏报。因此，有福的保禄自己也在这些话中大大赞美它：\BibleQuote{但如果我在肉身内生活，这是我工作的果实，那么我该选择什么，我不知道；}\BibleRef{斐 1:22}然后，将活在地上与离开并与基督同在的选项摆在自己面前，他选择了前者。但如果这生命是邪恶的，他就不会这样谈论它，也没有任何人，无论怎样努力，能把它引向德行的服务。因为没有人会用邪恶行善，用淫乱守贞，用嫉妒行仁爱。因此，当他说：\BibleQuote{肉性的思念不服天主的法律，也不能服，}\BibleRef{罗 8:7}他的意思是，恶本身不能成为德性；而\Quote{邪恶世界}这一表达必须理解为指邪恶的行为和败坏的道义原则。再者，基督来不是要处死我们，以此意使我们脱离现世生命，而是要把我们留在世上，预备我们相称地参与天上的住所。因此祂对圣父说：\BibleQuote{他们在世界上，我往祢那里去；我不求祢把他们从世界上拿去，但求祢保护他们脱免邪恶，}\BibleRef{若 17:11}即脱离罪恶。此外，那些不愿接受这一点，却坚持现世生命是邪恶的人，不应该责备自杀的人；因为正如那从邪恶中抽身的人不被责备，反而被认为配得冠冕，同样，那以暴力死亡、悬梁或其他方式结束自己生命的人，不应被定罪。然而，天主惩罚这样的人比惩罚杀人犯更重，我们都以恐惧看待他们，而且公正地如此行；因为如果毁灭他人是卑鄙的，那么毁灭自己就更卑鄙了。而且，如果这生命是邪恶的，杀人犯就应该得到冠冕，因为他们将我们从邪恶中解救出来。除此以外，他们被自己的话所困，因为他们将太阳置于他们神祇的第一等，月亮置于第二等，并敬拜它们为许多恩惠的赐予者，他们的说法是矛盾的。因为这些天体和其他天体的功用，无非是通过滋养和光照人畜的身体并使植物成熟，来服务于我们现世的生命——他们称之为邪恶的。那么，这\Quote{邪恶生命}的构造如何被那些按你们说是神祇的所如此服务呢？他们当然不是神，远非如此，而是天主为我们的使用所造的工程；这个世界也不是邪恶的。如果你对我说起杀人犯、奸淫者、盗墓者，这些事与现世生命无关，因为这些过犯不是来自我们在肉身内所过的生命，而是来自败坏的意志。因为，如果它们必然与这生命相连，如同与它同一命运，那么就没有人能免受它们，因为没有人能逃避人性的常态，诸如吃喝、睡眠、生长、饥渴、生老病死等等；没有人能超越这些，无论是罪人还是义人，君王还是农夫。我们都服从自然的必然性。所以，如果恶是这生命的本质要素，就没有人能避免它，正如刚才提到的事一样。不要对我说好人稀少，因为自然的必然性对所有人都不可超越，所以只要找到一个有德之人，我的论点就绝不会被削弱。可怜、可悲的人！你在说什么？这生命是邪恶的吗？在这生命中，我们得以认识天主，默想将来的事，成为天使而非人，参与天上万军的歌咏团？我们还需要什么证据来证明一个邪恶败坏的思想呢？\par

他们说：\Quote{那么，为什么}他们说，\Quote{保禄称现世生命为邪恶呢？}在称现世（时代）为邪恶时，他适应了我们的惯用法，我们常说：\Quote{我今天过得很糟}，这样抱怨的不是时间本身，而是行为或环境。同样，保禄在抱怨邪恶的行为原则时，使用了这些习惯用语；他表明基督既救我们脱离过犯，也为我们担保了将来。前者他已在以下话语中宣布：\BibleQuote{祂为我们自己的罪舍己；}又加上：\BibleQuote{为要救我们脱离这邪恶的现世，}他宣告了我们未来的安全。因为这两件事，法律都无能为力，但恩宠对二者都绰绰有余。\par

第4节。\BibleQuote{按照我们的天主和圣父的旨意。}\BibleRef{迦 1:4}\par

因为他们被这样的想法吓住了：如果离弃那旧法律而持守新法律，他们就会违背赐下法律的天主，所以他纠正他们的错误，说这事也悦纳了圣父：而且不单是\Quote{圣父}，而是\Quote{我们的圣父}，他这样做是为要感动他们，表明基督已使祂的圣父成为我们的圣父。\par

第5节。\BibleQuote{愿光荣归于祂，直到世世代代。阿们。}\BibleRef{迦 1:5}\par

这也是新颖而不寻常的，因为我们从未发现\Quote{阿们}一词放在书信的开头，而是在稍后的地方；然而，他在这里的开端就用它，表明他先前所说的话已包含了对迦拉达人的充分指控，并且他的论证是完整的，因为明显的过犯不需要详尽的控告。在宣讲了十字架、复活、从罪恶中的救赎和未来的保障、圣父的旨意和圣子的意愿、恩宠与平安以及祂整个恩赐之后，他以赞美结束。\par

另一个原因是他被恩赐的伟大、恩宠的丰盈、我们原本是谁以及天主所成就的事所产生的极大惊愕，而且是同时在一瞬间发生的。无法用言语表达这一点，他迸发出一篇赞美颂，为整个世界献上一首颂词，虽然确实不配这主题，却是他所能做到的。因此，他也开始使用更激烈的言辞；好像被对天主恩惠的感受所大大激动，因为他说了\Quote{愿光荣归于祂，直到永远。阿们。}之后，便开始更严厉的责备。\par

第六节。\BibleQuote{我惊奇你们竟这么快离开那以基督的恩宠召你们的那一位，去归向另一个福音。}\BibleRef{迦 1:6}\par

如同那些迫害基督的犹太人一样，他们想象自己遵守法律是为圣父所悦纳的，因此他指出，这样做不仅得罪了基督，也得罪了圣父，因为他们不但脱离了基督，也脱离了圣父。如同旧盟约不仅是由圣父赐予，也是由圣子赐予的，恩宠的盟约同样出自圣父与圣子，他们的一切作为都是共通的：\BibleQuote{凡父所有的都是我的。}\BibleRef{若 15:16} 通过说他们已从圣父那里坠落，他给他们带来双重指控：背信，以及立即的背信。相反的极端——迟来的背信——也是应受责备的，但一个人在第一波攻击中，就在初次交锋中坠落，显示的是极端懦弱的例子，而他也恰恰以这事来控告他们，说：\Quote{这是怎么回事呢？你们的诱惑者甚至不需要时间策划，第一次接触就足以使你们颠覆和俘获？你们有什么借口？如果这在朋友之间是一种罪行，抛弃旧日有用同伴的人应受谴责，那么，想想那背叛召叫他的天主的人应受何等惩罚！}他说：\Quote{我惊奇，}不仅是为了责备他们，在如此丰厚的恩赐、如此罪过的赦免、如此洋溢的仁慈之后，他们竟逃脱到奴役的轭下，而且也是为了表明他对他们所抱有的看法是良好而崇高的。因为他若将他们列在普通易受欺骗的人中，就不会感到惊奇。\Quote{但是既然你们，}他说，\Quote{是高贵的类型，并且受了许多苦，我实在惊奇。}这当然足以挽回他们，引导他们回到最初的表白。他在书信的中段也提到这一点：\BibleQuote{你们竟白白地受了那么多苦吗？如果真是白白地。}\BibleRef{迦 3:4} \Quote{你们正在离去；}他说的不是\Quote{你们已经离去，}也就是说，\Quote{我不相信或假定你们的诱惑已经完成；}这是一个即将挽回他们的人的语言，他在后面更清楚地用以下话语表达：\BibleQuote{我在主里深信你们不会有别的想法。}\BibleRef{迦 5:10}\par

\Quote{从那位以基督的恩宠召你们的那一位。}\par

蒙召来自圣父，但原因在于圣子。是祂带来了和好并作为恩赐赐予，因为我们不是因行为成义而得救；或者我更应该说，这些祝福来自两者；正如祂说：\BibleQuote{我的是你的，你的是我的。}\BibleRef{若 17:10} 祂说的不是\Quote{你们正在离开福音}，而是\Quote{离开那召你们的天主}，这是一种更可怕的说法，也更能影响他们。他们的诱惑者不是突然行事，而是逐步进行，实际上当他们将他们从信仰中移开时，却保留了名称不变。这是撒殚的策略，不公开设下陷阱；如果他们催促他们离开基督，他们会作为欺骗者和败坏者被避开，但容许他们继续停留在信仰中，并将他们的错误披上福音的名称，毫无畏惧地他们挖掘建筑物的基础，使用他们所用的术语作为掩盖毁灭者本身的帘幕。因此，既然他们给这骗局冠以福音的名，他就针对这名本身进行抗争，并大胆说：\Quote{归向另一个福音，}——\par

第七节。\BibleQuote{那并不是另一个福音。}\BibleRef{迦 1:7}\par

这说得合情合理，因为再没有别的福音。然而，\OriginalTerm{马尔基翁派}{Marcionites}被这句话误导，如同病人甚至被健康的食物伤害一样；因为他们抓住了这句话，并高喊：\Quote{保禄本人已宣称没有别的福音。}因为他们不承认所有的福音作者，只承认一位，而且按照他们的喜好，将福音书残缺不全、混乱不堪。他们对\BibleQuote{按照我的福音和耶稣基督的宣讲}这句话的解释，\BibleRef{罗 16:25}是极其可笑的；但为了那些容易受诱惑的人，有必要反驳它。因此，我们断言：即使写了一千部福音，如果所有内容都相同，它们仍然是一部福音，而它们的统一丝毫不会因作者的人数而受损。同样，另一方面，如果只有一个作者，但他却自相矛盾，那么所写内容的统一就会被破坏。因为一部作品的一致不在于作者的数量，而在于内容的一致或矛盾。由此可见，四部福音是一部福音；因为，既然四部福音都说同一件事，其统一就由内容的和谐得以保存，而不因人的不同而受损。保禄现在讲的不是数量，而是所言内容的差异。如果\BibleBook{玛窦}{福音}和\BibleBook{路加}{福音}在内容的意义和教义的准确性上有所不同，他们抓住这句话或许有理；但既然它们是一样的，那就让他们停止愚蠢，假装对连孩童都明白的事一无所知吧。\par

第7节：\BibleQuote{那扰乱你们、并想改变基督福音的，不过有些人罢了。} \BibleRef{迦 1:7}\par

这就是说，只要你们的理智清醒，只要你们的眼光健康，不受扭曲和虚幻的幻象所困，你们就不会承认另一个福音。因为正如有病的眼睛认错呈现在它面前的对象，同样，当心灵被恶念的混乱搅浑时，也是如此。因此，疯子混淆对象；但这种疯狂比身体的疾病更危险，因为它伤害的不是感官领域，而是心灵领域；它造成的混乱不是身体视觉器官的问题，而是理解之眼的问题。\par

\Quote{要更改基督的福音。}他们事实上只引进了两条诫命，即割损与守日子，但他说福音已被颠覆，以表明细微的掺假足以败坏全局。因为正如一个人只刮去王币上的部分肖像，就使整枚钱币变为伪币；同样，凡稍稍偏离纯正信德的人，很快会从这一点走向更严重的错误，并完全败坏。那么，那些指责我们好争论，因与异端分离，并说我们之间除了出于野心之外并无真正差异的人在哪里？让他们听听保禄的断言：那些稍有创新的人，就已经颠覆了福音。说圣子是受造之物，难道是一件小事？你们不知道吗？即使在旧盟约之下，一个人在安息日拾柴，违犯了一条诫命，而且不是大诫命，就被处死？\BibleRef{户 15:32}还有乌匝，当约柜快要倾倒时他扶住它，却突然被击毙，因为他侵入了不属他的职务？\BibleRef{撒下 6:6}因此，如果违犯安息日、触摸将要倾倒的约柜，如此显著地招致天主的义怒，甚至不给犯者片刻喘息，那么那败坏无比可怕的道理的人，难道还能找到借口和宽恕吗？肯定不会。在小事上缺少热心是一切灾祸的原因；因为轻微的错谬逃脱了适当的纠正，更大的错误便悄悄侵入。如同在身体上，忽略伤口会引起发烧、坏死和死亡；同样，在灵魂上，被忽视的小恶为更大的恶打开了门。一个人忽略守斋被看作小错；另一个人确立了纯正的信德，却因环境而装假，放弃他大胆的宣认，这也不算大事或可怕的事；第三个人被激怒，威胁要离开真信德，就以激情和怨恨为借口得到原谅。这样，每天有成千上万的类似错误被引入教会，我们成了犹太人和希腊人的笑柄，因为教会分裂成了上千个派别。但如果当初对那些企图轻微歪曲、偏离天主神谕的人给予了适当的斥责，就不会产生这样的瘟疫，这样的风暴也不会袭击众教会。现在你们会明白为什么保禄称割损是对福音的颠覆。如今我们中间有许多人，与犹太人同一天守斋，并以同样的方式守安息日；而我们容忍这件事，高尚地说，或者更恰当地说，是卑劣和可耻地容忍了。我为什么提及犹太人呢？因为看到我们中间有些人遵守许多外邦人的习俗：预兆、占卜、征兆、日子的区分、对子女出生情况的过分关注，以及他们一出生就将有不敬虔铭文的牌子放在他们不幸的头上，从而自始就教导他们放弃德行的努力，并至少将他们的一部分置于命运的虚假统治之下。但如果基督对受割损的人毫无益处，那么对于那些引入了这样败坏的人，信德将来对他们的救恩有何帮助呢？虽然割损是由天主赐下的，但保禄仍竭力废除它，因为不合时宜地遵守它有损于福音。那么，如果他对不当持守犹太习俗如此热切反对，我们对于不废除外邦习俗还有什么借口呢？因此，我们的事务如今陷入混乱和困境，因此我们的学生们充满了骄傲，颠倒了事物的秩序，使一切陷入混乱；而我们这些管理者忽视了对他们的管教，他们便鄙视我们哪怕温和的责备。即使他们的长上更加无用，充满无数邪恶，门徒也不应该违背。论到犹太经师，据说他们坐在梅瑟的位子上，门徒必须服从他们，虽然他们的行为如此邪恶，以致主禁止祂的门徒效法他们。因此，那些侮辱和践踏人——教会领袖、且借着天主的恩宠度圣洁生活的人——还有什么借口呢？如果我们彼此判断是不合法的，那么判断我们的老师就更不合法的。\par

第8-9节。\Quote{但是，即使是我们，或是一位从天上来的天使，给你们宣讲一个与我们以前所宣讲的不同福音，他就该受\OriginalTerm{绝罚}{anathema}。}\par

请看宗徒的智慧；为了避免有人指责他出于虚荣而赞扬自己的道理，他也把自己包括在\OriginalTerm{诅咒}{anathema}之中；由于他们转而求助于权威，即雅各伯和若望的权威，他也提到了天使，说：\Quote{请勿对我提及雅各伯和若望；若天上最崇高的天使之一败坏了福音，就让他受诅咒。}\Quote{天上}这个词是有意加上的，因为司铎也被称为天使。\BibleQuote{因为司铎的唇舌应保存知识，人们应从他的口中寻求法律，因为他是万军之上主的使者。}\BibleRef{拉 2:7}因此，为了避免人们以为这里的天使是指司铎，他通过加上\Quote{从天上}指明了天上的灵体。他说，不是如果他们宣讲相反的福音，或颠覆了全部真正的福音，就让他们受诅咒；而是如果他们稍微改变，或偶然扰乱了我的道理。\Quote{正如我们从前说过，如今我再说：}为使他的话不显得出于愤怒、夸张或鲁莽，他如今重复了这些话。情感也许会因愤怒而改变，但重复第二次证明这话是经过深思熟虑的，并且先前已得到判断的认可。当亚巴郎被请求派遣拉匝禄时，他回答说：\BibleQuote{他们自有梅瑟和先知，听从他们好了。如果他们不听从他们，纵使有人从死者中复活，他们也不会信服。}\BibleRef{路 16:31}基督如此引述亚巴郎的话，是要表明他认为圣经比一个从死者中复活的人更值得信赖；保禄也是如此（我说保禄，是指引导他思想的基督），宁可信圣经，也不信从天上降下来的天使。这是合理的，因为天使虽然强大，但只是仆人和仆役，而圣经不是由仆人，而是由万有的天主所写成和派遣的。他说，若\Quote{任何人}向你们宣讲与我们先前所宣讲的不同福音——而非\Quote{若这人或那人}；在此显出了他的谨慎和避免得罪人的用心，因为当他在天上地下如此广泛地包括了所有一切时，何须再提名字呢？他诅咒福音传道者和天使，包括了每一种尊位，而他提到自己，包括了每一种亲密和关系。他喊道：\Quote{不要对我说，我的同宗徒和同僚曾如此说；若我宣讲这样的道理，我也不饶恕自己。}他说这话，并非谴责宗徒们偏离了他们受命传达的信息；绝无此意（因为他说，无论是我们还是他们如此宣讲）；而是要表明，在真理的讨论中，不应考虑人的身份尊卑。\par

第十节\BibleRef{迦 1:10}：\Quote{我现在是在说服人，还是在说服天主？或是我在寻求取悦人？若我仍取悦人，我就不是基督的仆人了。}\par

设若，他说，我可能用这些道理欺骗你们，我能欺骗天主吗？他知道我尚未说出的思想，而取悦他是我不断的努力。请看宗徒的精神，福音的崇高！同样，他写信给格林多人说：\BibleQuote{我们并不是再度向你们推荐自己，而是给你们机会夸耀；}\BibleRef{格后 5:12}又说：\BibleQuote{至于我，受你们审断，或受人的审断，为我是一件极小的事。}\BibleRef{格前 4:3}因为他不得不向弟子们解释自己，作为他们的导师，他屈服了；但他为之伤心，并非出于怨恨，绝非如此，而是因为那些被误导者的心思不稳定，以及未能得到他们完全的信任。因此，保禄现在似乎这样说：—— 我需向你们交代吗？我需受人的审判吗？我的账要算给天主，我的一切行为都着眼于那审判，我绝不会如此可怜地堕落去歪曲我的道理，因为我所宣讲的要在万有的主面前获得辩白。\par

他这样说，既是为反驳他们的意见，也是为自己辩护；因为门徒的本分是服从师傅，而非判断师傅。但现在，他说，次序颠倒了，你们竟坐在审判席上，要知道我几乎不在乎在你们面前为自己辩护；我所做的一切都是为了天主，好能向祂交代我的道理。凡想说服人的，就会行事迂回不真诚，并运用欺骗和虚假，以谋取听众的赞同。但凡是向天主说话、渴望取悦祂的人，则需要心灵的纯朴与洁净，因为天主不能被欺骗。由此可见，我给你们写信并非出于统治欲，或为招收门徒，或为从你们手中得到荣誉。我的努力一直是取悦天主，而非取悦人。否则，我仍该与犹太人来往，仍该迫害教会——我完全抛弃了祖国、同伴、朋友、亲属和一切声望，换来了迫害、仇恨、争斗和每日逼近的死亡，这便是一个明显的证据，证明我不是出于爱人的称赞而说话。他说这话，是即将叙述他先前的生平、突然的皈依，并清楚证明这是真诚的。为使门徒不致因以为他是在为自己辩白而自高自大，他先说了：\Quote{我如今是劝人吗？} 他很清楚在适当的时候，以庄重而高昂的语调纠正门徒；当然，他还有其他证据来证明他宣讲的真理性——通过神迹异事、危险、监禁、每日的死亡、饥渴、赤身露体等。但现在他谈论的不是假宗徒，而是分享过这些危险的真宗徒，他用了另一种方法。因为当他指向假宗徒时，他通过提出自己所忍受的危险来比较，说：\BibleQuote{他们是基督的仆人吗？（我说话如同疯狂）我更是；劳苦更甚，监禁更甚，鞭打远重，屡次冒死。} \BibleRef{格后 11:23} 但现在他谈论他先前的行事方式，并说：\par

第11、12节。\BibleQuote{弟兄们，我愿意你们知道，我所传的福音并不是出于人。因为我既不是从人领受的，也不是被教导的，而是借耶稣基督的启示而来的。}\par

你们看，他是多么谨慎地断言自己由基督所教导，基督亲自无需人的介入，屈尊俯就向他启示了一切知识。若有人问他有何凭据，证明天主就这样直接启示了这些不可言喻的奥秘，他就会举出他先前的行事方式，论证他的皈依若非出于天主的启示，不会如此突然。因为当人激烈而热切地倾向相反的一面时，他们的归信若由人的方法促成，则需要大量的时间和技巧。因此，显然那突然皈依，并在疯狂至极时清醒过来的人，必定蒙受了天主的启示和教导，从而立刻获得了完全的神智。为此，他不得不叙述他先前的生平，并召请迦拉达人为过去的事作证。他说：天主独生子亲自从天赏赐召叫了我，你们这些不在场的人无法知道，但你们确实知道我从前是个迫害者。因为我的暴行甚至传到你们耳中，而巴勒斯坦与迦拉达之间的距离如此之大，若非我的行为超越了所有界限和忍耐，这消息不会传到那里。因此他说：\par

第13节。\BibleQuote{你们听说过我从前在犹太教中的行事方式，我是怎样过度地迫害天主的教会，并极力摧残她。} \BibleRef{迦 1:13}\par

请注意，他毫不回避加重每个要点；他不仅简单地说他\Quote{迫害}，而且说\Quote{过度地}，不仅说\Quote{迫害}，而且说\Quote{极力摧残}，这表示一种试图消灭、拆毁、破坏、毁灭教会的企图。\par

第14节。\BibleQuote{我在犹太教中，比我本国许多同辈的人更有长进，对祖传的传统更为热衷。} \BibleRef{迦 1:14}\par

为避免造成他的迫害出于情感、虚荣或敌意的印象，他表明自己是受热忱所驱使，固然不是\Quote{按着知识}，\BibleRef{罗 10:2} 但仍是对祖传传统的一种热切推崇。他的论证如下：如果我对教会的攻击不是出于人的动机，而是出于宗教的却是错误的热忱，那么如今我为教会奋斗并接受了真理，我又怎会被虚荣心所驱使呢？如果在我错误时，支配我的不是这种动机，而是虔诚的热忱，那么如今我认识了真理，就更不该受这种猜疑了。我一转入教会教义，就摈弃了犹太人的偏见，在那一方面表现出更加炽烈的热忱；这证明我的皈依是真诚的，而且支配我的热忱是来自天上。还有什么别的理由能促使我作出这样的改变，拿荣誉去换轻蔑，拿安息去换危险，拿安全去换困苦？除了对真理的爱，绝无其他。\par

第15、16节。\BibleQuote{然而，当那从我母腹中分别了我，并借祂的恩宠召叫我的天主，乐意将祂的圣子启示在我里面，叫我把祂传扬在外邦人当中，我就没有与血肉商量。}\par

他的目的在此是要显示，他是因着某种隐秘的眷顾，暂且任由他自己。因为若他从母腹中被分别出来，为要作宗徒并被召叫去从事那职事，却直到那时才实际被召，而他一听到召唤便立即服从，那么显然天主对这段迟延必有深意。你们或许急于从我这里得知这目的是什么，尤其是为何他没有与十二位一同被召。但为免因离题去讨论更紧迫的事而拖延这篇讲辞，我必须恳请你们的爱德，不要向我索求一切，而要自行搜寻，并向天主祈求，求祂启示给你们。此外，在我之前向你们论述他名字由扫罗改为保禄的讲道中，我曾部分讨论过这主题；你们若已忘记，从阅读那卷书便可完全明了。眼前让我们继续我们讲论的线索，并思考他现在所提出的证据，证明所发生的事并非出于自然——乃是天主亲自以眷顾安排了此事。\par

\BibleQuote{因祂的恩宠召叫我。}\par

天主的确说，祂是因他杰出的能力而召叫他，正如祂对阿纳尼雅所说：\BibleQuote{因他是我所拣选的器皿，为把我的名字带到外邦人和君王面前}\BibleRef{宗 9:15}，也就是说，他堪当服务，并能成就大事。天主以此作为召叫他的理由。但他自己却处处将此归之于恩宠和天主不可言喻的仁慈，如那句话：\BibleQuote{然而我蒙受了怜悯}，并非因我堪当或甚至有用，而是\BibleQuote{为叫我在我身上显明耶稣基督所有的忍耐，给后来信祂得永生的人作榜样}\BibleRef{弟前 1:16}。看哪！他满溢的谦逊：我蒙受了怜悯，他说，为叫无人绝望，即使最坏的人也曾分享祂的恩惠。因为这话的意思正是：\BibleQuote{为叫祂显明祂所有的忍耐，给后来信祂的人作榜样。}\par

\BibleQuote{将祂的圣子启示在我里面。}\par

基督在别处说：\BibleQuote{除了父，没有人知道子是谁；除了子，和子所愿意启示的人，也没有人知道父是谁}\BibleRef{路 10:22}。您们看到，父启示子，子也启示父；同样，论到祂们的荣耀，子光荣父，父也光荣子；\BibleQuote{光荣您的子，使子也光荣您}，以及\BibleQuote{如同我光荣了您}\BibleRef{若 17:1}。但祂为什么说\BibleQuote{将祂的圣子启示在我里面}，而不说\Quote{启示给我}？\par

\BibleQuote{为叫我在外邦人中传扬祂。}因为不仅他的信德，连他蒙选召领受宗徒职分，也出于天主。他说，祂这样特别启示给我，不只是为叫我亲自看见祂，也为叫我将祂显明给别人。他说的不仅是\Quote{别人}，而是\BibleQuote{为叫我在外邦人中传扬祂}，如此预先触及他那辩护的重要根据，即在于门徒各自的特征；因为向犹太人和向外邦人宣讲的方式是必须不同的。\par

\BibleQuote{我并没有立即与血肉之人商量。}\par

这里他暗指宗徒们，以他们的肉体本性称呼他们；不过，我也不会否认他可能指全人类。\par

17节。\BibleQuote{也没有上耶路撒冷去见那些比我先作宗徒的人。}\BibleRef{迦 1:17}\par

这些话语孤立地看，似乎散发出一种傲慢的精神，与宗徒的气质格格不入。因为为自己投票，不允许任何人分享自己的意见，是愚妄的标志。经上记载：\Quote{\BibleQuote{人若自以为是，愚人比他更有希望；}}\BibleRef{箴 26:12}又说：\Quote{\BibleQuote{那些自视为聪明、自以为精明的人，有祸了！}}\BibleRef{依 5:21}而\Person{保禄}自己在另一处也说：\Quote{\BibleQuote{不要自以为聪明。}}\BibleRef{罗 12:16}受过这样教导、又这样劝戒他人的人，即使是个普通人，也不会陷入这样的错误；更何况是\Person{保禄}本人呢？然而，正如我所说，如果单看这些话，很容易成为许多听众的陷阱和冒犯。但如果把原因附加进去，所有人都会称赞和钦佩说话的人。那么，我们就这样做吧；因为正确的做法不是衡量字句，也不是单独检查语言，否则会导致许多错误，而是要关注作者的意图。如果我们不在自己的言论中遵循这个方法，不去探究说话者的心意，我们就会树敌众多，一切都会陷入混乱。这不仅限于言语，如果不遵守这个规则，行动也会产生同样的结果。因为外科医生常常切割和折断某些骨头，强盗也这样做；然而，如果不能区分二者，那将是可悲的。同样，杀人者和殉道者受折磨时经历同样的痛苦，但两者之间的区别很大。如果我们不注意这个规则，我们将在这些事情上无法辨别；我们会把\Person{厄里亚}、\Person{撒慕尔}和\Person{丕乃哈斯}称为杀人者，把\Person{亚巴郎}称为杀子者；也就是说，如果我们只审视表面的事实，而不考虑行为者的意图。那么，让我们来探究\Person{保禄}这样写的意图，让我们考虑他的目的和对宗徒们的总体态度，以便我们理解他现时的意思。他过去和现在都不是为了贬低宗徒或抬高自己而说话（这怎么可能？他把自己也包括在他的绝罚之下？），而总是为了维护福音的纯正。因为扰乱教会的人说，应该服从遵守这些规条的宗徒，而不是禁止它们的\Person{保禄}，因此犹太主义异端逐渐渗入，所以他必须勇敢地抵抗他们，出于压制那些妄自尊大的人的傲慢，而不是为了说宗徒的坏话。因此，他说：\Quote{我没有与血肉商议；}因为一个受天主教导的人，后来又去请教人，那是极其荒谬的。因为向人学习的人，反过来应该把人当作自己的顾问。但那位蒙受了神圣而蒙福的声音，并且由拥有全部智慧宝藏的那位完全教导的人，为何还要去与人商议呢？他应当教导，而不是被人教导。因此，他这样说，不是出于傲慢，而是为了显示自己使命的尊严。他说：\Quote{我也没有上耶路撒冷去见那些在我以前作宗徒的。}因为他们不断重复说宗徒在他以前，且在他以前蒙召，所以他说：\Quote{我没有上他们那里去。}如果他有必要与他们交通，那启示他使命的那位定会给他这样的命令。然而，他真的没有上那里去吗？不，他上去了，而且不仅是为了这个，还是为了向他们学习一些东西。当关于我们当前主题的问题在\Place{安提约基雅}城中，在起初就如此热心的教会中产生时，并且讨论了外邦信徒是否该受割损，或者没有必要接受这个仪式时，正是\Person{保禄}本人和\Person{息拉}上去了。那么，他怎么说他没有上去，也没有商议呢？首先，因为他不是主动上去的，而是被他人派遣的；其次，因为他不是去学习，而是去说服别人。因为他从一开始就持这样的见解，后来宗徒们批准了，即割损是不必要的。但当这些人认为他不值得信任，而向\Place{耶路撒冷}的人求助时，他上去不是要进一步受教，而是要说服反对者，让\Place{耶路撒冷}的人与他一致。因此，他从一开始就看到了正确的行动路线，不需要老师，而是在任何讨论之前，就坚定不移地坚持了宗徒们经过大量讨论\BibleRef{宗 15:2}后来批准的做法。\Person{路加}在他的记述中表明，\Person{保禄}在去\Place{耶路撒冷}之前曾与他们就此问题进行了长时间的辩论。但由于弟兄们选择通过\Place{耶路撒冷}的人了解这件事，他是为他们而不是为自己上去的。他说\Quote{我没有上去}这句话，表示他既不是在教导之初上去的，也不是为了受教的目的。两者都包含在\Quote{我立即没有与血肉商议}这句话中。他不只说\Quote{我没有商议}，而是说\Quote{立即}；而他后来的行程也不是为了获得更多教导。\par

\Quote{\BibleQuote{我却往阿拉伯去了。}}\BibleRef{迦 1:17}\par

看，一个热忱的灵魂！他渴望占据尚未开垦、仍处荒芜的地区。若他与宗徒们同住，既无可学习，他的宣讲就会受到限制，因为他们必须将圣言传播到各处。因此，这位蒙福者，心灵热忱，立即着手教导野蛮的异邦人，选择了充满战斗与劳苦的生活。说了\Quote{我去了阿拉伯}后，他又说\Quote{然后我又回到了大马士革}。在此请观察他的谦卑：他未提及自己的成功，也未提及他教导了谁、有多少人。然而，他的热忱在圣洗后立即如此强烈，以致他驳倒了犹太人，激怒了他们，以至于犹太人和希腊人设计谋害他。若不是他大大增加了信徒的人数，这不会发生；因为他们在教理上被击败，便诉诸谋杀，这明显是\Person{保禄}优越的标记。但基督不容他被害，为他保存了使命。然而，他对这些成功只字不提，因此在他所有的言论中，他的动机并非野心，也不是要获得比宗徒更高的尊荣，也不是因被轻视而懊恼，而是担心他的使命会受到任何损害。因为他称自己为\Quote{未及产期而生的人}、\Quote{罪人中的罪魁}、\Quote{宗徒中最小的}、\Quote{不配称为宗徒}。说这话的人，是劳作比众人更多的；这才是真正的谦卑：一个自觉毫无卓越之处而谦卑自述的人，是诚实而非谦卑；但在取得如此胜利后仍如此说，则是操练节制。\par

第17节。\BibleQuote{然后我又回到了大马士革。}\BibleRef{迦 1:17}\par

然而，他在此城岂非成就了何等大事？因为他告诉我们，\Person{阿勒达}王手下的总督派人围住全城，企图捉拿这位蒙福者。这有力地证明他受到犹太人猛烈的迫害。然而，在此他对这些只字未提，只提到他的到达与离开，对那里发生的事保持沉默；若非情况需要叙述，他在我提到的那个地方也不会提及这事。\BibleRef{格后 11:32}\par

第18节。\BibleQuote{此后三年，我上耶路撒冷去拜访\Person{刻法}。}\BibleRef{迦 1:18}\par

有什么比这样的灵魂更为谦卑呢？在如此成功之后，他并不需要\Person{伯多禄}任何东西，甚至不需要他的认可，尽管与他有同等的尊荣（目前我不再多说），却仍以他为长上和上级而来见他。此次旅程的唯一目的就是拜访\Person{伯多禄}；因此，他向宗徒们表达应有的尊敬，并估量自己不但不是他们的上辈，甚至不是他们的平等。这从此次旅程可见一斑：\Person{保禄}受感动去拜访\Person{伯多禄}，正是出于许多弟兄与圣人同住时的同样心情；或者更谦卑地说，因为他们这样做是为了自己的益处，而这位蒙福者不是为了自己的受教或改正，仅仅是为了因他的临在而瞻仰和尊崇\Person{伯多禄}。他说\Quote{拜访\Person{伯多禄}}；他说的不是\Quote{看}（\OriginalTerm{看}{\GreekText{ἰδεῖν}}），而是\Quote{拜访并察看}（\OriginalTerm{拜访并察看}{\GreekText{ἰστορῆσαι}}），那些想认识伟大而辉煌之城的人常用此词。他认为\Person{伯多禄}的面容本身值得如此辛劳；这一点也见于\WorkTitle{宗徒大事录}\BibleRef{宗 21:17}。因为另一次到\Place{耶路撒冷}，在使许多外邦人归化，并以远超他人的劳苦将\Place{旁非利亚}、\Place{吕考尼亚}、\Place{基里基雅}以及世界那一角的所有民族改造并带到基督面前之后，他首先极其谦卑地向\Person{雅各伯}说话，视他为长上和上级。接着他听从了他的劝告，而那劝告与这封信相反：\Quote{弟兄，你看犹太人中有多少相信的人，成千上万；因此你剃头，洁净自己。}\BibleRef{宗 21:20} 于是他剃了头，遵守了所有犹太礼仪；因为凡\WorkTitle{福音}不受影响的地方，他是众人中最谦卑的。但若他看出因这种谦卑而有人受损害，他就放弃这种不当的谦卑，因为那不再是谦卑，而是侮辱和毁灭门徒。\par

第18节。\BibleQuote{与他同住了十五天。}\BibleRef{迦 1:18}\par

为他而旅行是尊敬的标志；但逗留这么多天，则是友谊与最诚挚的爱的标志。\par

第19节。\BibleQuote{至于其他的宗徒，除了主的兄弟\Person{雅各伯}，我没有见过一位。}\BibleRef{迦 1:19}\par

请看，他与\Person{伯多禄}的友谊何等深厚！为了他，他离开家乡，与他同住。我屡次提及此事，并愿你们记住，以免任何人听到这位宗徒似乎曾反对\Person{伯多禄}的话时，对他产生怀疑。他事先声明，当他说\Quote{我反对了\Person{伯多禄}}时，无人会以为这些话隐含敌意与争辩；因为他尊崇并爱戴\Person{伯多禄}本人胜于一切，此次旅程仅为他一人，并非为其他任何人。\Quote{至于其他的宗徒，除了\Person{雅各伯}，我没有见过一位。}他的意思是：\Quote{我仅仅见到了他，并非从他学习。}但请看，他多么荣幸地提及他：他不单说\Quote{\Person{雅各伯}}，还加上这个显赫的称号，显示他毫无嫉妒。若他仅想指明所指何人，他可以用另一个称呼，称他为\Person{克罗帕}的儿子，如同圣史所作的那样。但他认为\Person{雅各伯}分享宗徒的尊号，便通过尊崇\Person{雅各伯}来高抬自己；而这就是他称他为\Quote{主的兄弟}的原因，尽管\Person{雅各伯}并非按血统是主的兄弟，只是这样被称为。然而这并未阻止他赐予此称号；在许多其他事例中，他显露出对众宗徒的那种相称的高贵情操。\par

第20节。\BibleQuote{关于我写给你们的这些事，看，在天主面前，我不说谎。}\BibleRef{迦 1:20}\par

请整体观察这圣洁灵魂的透明谦卑：他为自己辩护的热切程度，犹如必须报告自己的行为，并在法庭上为生命辩护一样。\par

第21节。\BibleQuote{随后我来到了叙利亚和基里基雅地区。}\BibleRef{迦 1:21}\par

他谒见了伯多禄之后，便重新开始宣讲，并承担他面前的任务，他避开犹太地区，既因为他的使命是向外邦人传教，也因为他不愿意\Quote{建造在别人的根基上}。因此，从下文可以看出，甚至连偶然的相遇都没有。\par

第22、23节。\BibleQuote{我仍然以面容不被犹太的各教会所认识；他们只是听说：那从前迫害我们的，如今正宣讲他从前所摧残的信德。}\par

这是何等的谦逊！他再次提起自己迫害和摧残教会的事实，由此使他从前的生活蒙受恶名，却略过他即将成就的光辉事迹。他本可以讲述自己的一切成功，只要他愿意，但他一件也不提，而是用一个词跨过广阔的空间，仅仅说：\Quote{我来到了叙利亚和基里基雅地区；}以及\Quote{他们听说，那从前迫害我们的，如今正宣讲他从前所摧残的信德。} \Quote{我未被犹太的各教会所认识}这句话的目的，是要表明，他非但没有向他们宣讲割礼的必要性，甚至连面容也不曾为他们所知。\par

第24节。\BibleQuote{他们便因我而光荣了天主。}\BibleRef{迦 1:24} 这里再次看到，他是多么精确地遵守谦逊的准则；他并没有说：他们钦佩我、称赞我或对我感到惊奇，而是藉着这些话将一切归功于天主的恩宠：\BibleQuote{他们便因我而光荣了天主。}\par

\SourceRecord{Translated by Gross Alexander.；Nicene and Post-Nicene Fathers, First Series；Vol. 13.；Edited by Philip Schaff.；Buffalo, NY: Christian Literature Publishing Co.,；1889.；<http://www.newadvent.org/fathers/23101.htm>.}

% structure: p0001 source=h1 target=section reason=page-title

\section{迦拉达书讲道集 第二讲}

第1-2节\par

\Quote{\Emph{过了十四年}，我同\Person{巴尔纳伯}再上\Place{耶路撒冷}去，也带了\Person{弟铎}同去。\Emph{我是因着启示上去的。}}\par

他的第一次行程是出于渴望拜访\Person{伯多禄}，第二次，他说，是因\Person{圣神}的启示。\par

\BibleQuote{我将我在外邦人中所宣讲的福音向他们陈述，但在那些有权威的人前，我私下陈述，免得我白白地奔跑，或曾经空跑过。}\BibleRef{迦 2:2}\par

这是怎么回事，啊，保禄！你起初和三年后都没有与宗徒们商议，如今过了十四年，反倒与他们商议，免得你白白地奔跑？若起初就这样做，比过了这么多年更好；若你不确知自己没有空跑，又为何奔跑呢？谁会愚蠢到宣讲这么多年，而不确知自己的宣讲是真的？更增加难处的是，他说他是因启示上去的；然而，这个难处却为前一个难处提供了解决。若是他自愿上去，那实在极不合理，而且这有福的灵魂也不可能陷入这样的愚昧；因为正是他自己说：\BibleQuote{所以我奔跑，不像无定向的；我这样打拳，不像打空气的。}\BibleRef{格前 9:26} 因此，如果他奔跑\BibleQuote{不像无定向的}，又怎能说\Quote{免得我奔跑，或曾经空跑过？}呢？由此可见，如果没有启示就上去，他会做出愚昧的事。但实际情况并没有这样的荒谬；谁敢再存这种怀疑，既然是圣神的恩宠引他上去？为此他加了\Quote{因启示}一词，免得在问题解决之前，他被人判为愚昧；他深知这不是人的事，而是关乎现在和将来的深沉的天主眷顾。那么，他这次行程的原因是什么？正如他以前从\Place{安提约基雅}上\Place{耶路撒冷}去，不是为了自己（因为他清楚知道，他的责任只是服从基督的教训），而是为了说服那些好争辩的人；同样，现在他的目的是使控告者完全满意，并非他自己想知道自己没有空跑。他们认为\Person{伯多禄}和\Person{若望}（他们对他们的评价比保禄高）与他不同，因为他在宣讲中省略了割损礼，而前者却允许它，因此他们相信他这样做是违法的，是在空跑。他说，我上去，将我的福音传达给他们，并非我自己要学习什么（从后面可以更清楚地看出），而是要使这些怀疑的人确信我没有空跑。圣神预见了这场争论，便安排他上去做这次传达。\par

因此他说他是因启示上去的，并带着\Person{巴尔纳伯}和\Person{弟铎}作为他宣讲的见证人，将自己在外邦人中宣讲的福音（即省略了割损礼）传达给他们。\Quote{但在那些有权威的人前，我私下陈述。}\Quote{私下}是什么意思？通常，一个人若想改革公议的教义，会公开提出，而非私下；但保禄这样做是为了私下的目的，不是要学习或改革什么，而是要堵住那些想要欺骗的人的口。在\Place{耶路撒冷}，若有人违反法律或禁止割损礼，所有人都感到愤慨；正如\Person{雅各伯}所说：\BibleQuote{兄弟，你看，在犹太人中信主的有多少千千万万，他们都听说你教训众犹太人背弃梅瑟。}\BibleRef{宗 21:20} 既然他们感到愤慨，他没有屈尊公开宣布他的宣讲是什么，而是在\Person{巴尔纳伯}和\Person{弟铎}面前与那些有权威的人私下商谈，以便他们能可信地向控告者作证：宗徒们在他的宣讲中没有发现不一致，反而予以确认。\Quote{那些有权威的人}（\OriginalTerm{那些有权威的人}{\GreekText{τοῖς} \GreekText{δοκοῦσιν}}）并不否认他们伟大的真实性；因为他论及自己时说：\Quote{我也\OriginalTerm{似乎}{\GreekText{δοκῶ}}有天主的神}，这并非否认事实，而是谦逊地陈述。这里的措辞也表达了他对共同意见的赞同。\par

\BibleQuote{但是，同我一起的\Person{弟铎}，虽是希腊人，也没有被强迫受割损。}\BibleRef{迦 2:3}\par

\Quote{虽是希腊人？}是什么意思？是说他有希腊血统，未受割损；因为我不但这样宣讲，\Person{弟铎}也这样做了，宗徒们并没有强迫他受割损。这清楚地证明，宗徒们并没有谴责保禄的教义或行为。不仅如此，甚至敌对一方对事实的紧迫陈述，也没有迫使宗徒们命令割损，正如他亲口所说——\par

\BibleQuote{这是因为那些偷偷进来的假弟兄。}\BibleRef{迦 2:4}\par

这里出现了一个非常重要的问题：这些假弟兄们是谁？如果宗徒们在耶路撒冷准许了割损，那么那些依照宗徒的判决而吩咐行割损的人，为何被称为假弟兄呢？首先，因为命令做一件事与在事后允许它是有区别的。命令者带着必行和首要的热忱去行；但自己不下令，却允许愿意者去做的人，并非出于必要，而是为了某种目的而让步。我们有一个类似的例子，在\WorkTitle{格林多前书}，他命令丈夫与妻子再同居。为避免被认为是在为他们立法，他随即补充说：\Quote{我说这话，是出于宽容，并不是出于命令。} \BibleRef{格前 7:5} 因为这并非权威性的判决，而是对他们纵欲的宽容；如他所说：\Quote{因你们的纵欲。} 你愿意知道保禄对此的看法吗？听他说：\Quote{我愿众人如同我自己一样，} \BibleRef{格前 7:7} 在节操方面。同样地，宗徒们在此让步，并非辩护律法，而是迁就犹太主义的软弱。如果他们辩护律法，就不会对犹太人用一种宣讲，对外邦人却用另一种了。如果割损对不信者是必要的，那么它显然对全体信徒同样也是必要的。但他们决定不在这一点上困扰外邦人，表明他们允许割损是出于对犹太人的迁就。而假弟兄们的目的是使他们脱离恩宠，再度被奴役。这是第一个区别，而且是很大的区别。第二个区别是，宗徒们只在犹大行动，那里律法仍然有效；但假弟兄们却处处行动，因为所有的迦拉达人都是受他们影响的。由此可见，他们的意图不是建立，而是完全拆毁福音；此事宗徒们是基于一个理由而容许，而假弟兄们则是基于另一个理由而热切实行。\par

第4节：\Quote{他们潜入，窥探我们在基督耶稣内的自由，为使我们再受束缚。} \BibleRef{迦 2:4}\par

他称他们为\Quote{窥探者}，以此指出他们的敌意；因为窥探者的唯一目的，就是通过了解对手的位置，为自己获得破坏和毁灭的便利。而这正是那些想要使门徒们回到旧日奴役的人所做的。由此也可见，他们的目的与宗徒们是何等截然相反；后者做出让步是为了逐渐使他们脱离奴役，而前者则密谋使他们遭受更严重的奴役。因此，他们四处窥探，仔细观察，多管闲事，找出哪些人未受割损；正如保禄所说：\Quote{他们潜入，窥探我们的自由，} 这不仅用\Quote{窥探者}一词，还用这种潜入和偷偷摸摸的表达，指出他们的阴谋。\par

第5节：\Quote{我们就是一刻也没有顺从屈服于他们。} \BibleRef{迦 2:5}\par

请注意这句话的力量和强调；他说的不是\Quote{通过辩论}，而是\Quote{通过屈服}，因为他们的目的不是教导正确的教义，而是征服和奴役他们。因此，他说，我们向宗徒们让步，但不向这些人让步。\par

第5节：\Quote{为叫福音的真理存留在你们中间。} \BibleRef{迦 2:5}\par

他说，这是为了使我们通过行动证实我们已经用言语所宣告的——也就是说，\Quote{旧事已过，看，都成了新的；} 和 \Quote{谁若在基督内，他就是一个新受造物；} \BibleRef{格后 5:17} 以及 \Quote{若你们接受割损，基督就对你们毫无益处。} \BibleRef{迦 5:2} 为了持守这个真理，我们一刻也没有让步。然后，由于宗徒们的行为直接与他相悖，而且可能有人会问他们命令行割损的理由，他便着手解决这个异议。他很有技巧地处理，因为他没有给出真正的理由，即宗徒们的行动是出于迁就和一种安排（\OriginalTerm{经纶}{\GreekText{οἰκονομία}}）似的；否则听众就会受到损害。因为那些要从安排中获益的人，不应该知道它的设计；一旦暴露，一切就都毁了。因此，参与安排的人应该知道其意图，而受益的人则不应该知道。为了更清楚地说明我的意思，我从当前的主题中举一个例子：蒙福的保禄本人，他本意要废除割损，但当他要派遣弟茂德去教导犹太人时，却先给他行了割损，然后才派他去。他这样做，是为了使听众更容易接纳他；他以割损开始，为的是最终废除它。但这个理由他只告诉了弟茂德，而没有告诉门徒们。如果他们知道他行割损的目的正是为了废除这仪式，他们就永远不会听他的宣讲，整个益处也就丧失了。但现在，他们的无知对他们最有益，因为他们认为他的行为是出于对律法的尊重，这导致他们以善意和礼貌接纳了他和他的教导，并且逐渐接纳他并受到教导后，他们就放弃了旧习俗。如果他们起初就知道他的理由，就不会发生这种情况；因为他们会避开他，并且避开后就不会听他讲，不听讲就会继续留在以前的错误中。为了防止这一点，他没有透露他的理由；在这里他也没有解释这安排（\OriginalTerm{经纶}{\GreekText{οἰκονομία}}）的缘由，而是以不同的方式组织他的言论，如此：\par

第6节：\Quote{至于那些被视为重要的人（不论他们以前是何等样人，与我毫不相干；天主不以外貌取人。）} \BibleRef{迦 2:6}\par

此處他不僅不為宗徒們辯護，反而嚴厲指責那些聖人，是為了軟弱者的益處。他的意思是：雖然他們准許割損，但他們將向天主交賬，因為天主不看重他們的身分地位。然而他沒有說得那麼直白，而是謹慎地表述。他沒有說：如果他們敗壞了教義，偏離了傳道所定的規範，必受最嚴厲的審判並受罰；而是更恭敬地暗示他們，用這些話：\Quote{那些被認為是有分量的，無論他們從前如何}。他沒有說\Quote{無論他們\InnerQuote{是}}，而是說\Quote{從前是}，表明他們從那時起也不再這樣宣講了，因為教義已普遍傳開。\Quote{無論他們從前如何}這句話，意指如果他們這樣宣講，他們就要交賬，因為他們要在天主面前，而非在人面前為自己辯白。他說這些，並非對他們行為的正確性有所懷疑或不知，而是（如我所說）出於對這樣組織言論的適宜性的考慮。接著，為了不讓人覺得他站在對立面指控他們，從而產生他們意見不合的猜疑，他立即加上這一修正：\Quote{那些被認為是有分量的，在會議中並沒有加給我什麼}。他的意思是：你們可能怎麼說，我不知道；我很清楚的是，宗徒們並沒有反對我，而是我們的意見一致、協調。這從他的措辭\Quote{他們給了我共融的右手}可以看出；但他此刻還沒有說這句話，只是說他們沒有在任何方面告知或糾正他，也沒有增加他的知識。\par

第6節。\BibleQuote{那些被認為是有分量的，並沒有加給我什麼：}\BibleRef{迦 2:6}\par

這就是說，當他們得知我的行動時，他們沒有添加什麼，也沒有糾正什麼；儘管他們知道我的旅程目的就是要與他們共融，我是因聖神的啟示而來，並且我帶著未受割損的弟鐸同來，他們既沒有給他行割損，也沒有給我增加任何額外的知識。\par

第7節。\BibleQuote{反倒。}\BibleRef{迦 2:7}\par

有人認為他的意思是：宗徒們不僅沒有教導他，反而受他教導。但我不會這樣說，因為他們每一位都已完全受過教導，又能從他那裡學到什麼呢？因此他並不是通過\Quote{反倒}這個表達來意指這一點，而是說他們不僅沒有責備他，反而讚揚他：因為讚揚是責備的相反。有人可能會在此回應：如果宗徒們讚揚你的做法，那麼按照邏輯，他們就該廢除割損？現在，保祿認為斷言他們廢除了割損過於大膽，也與他自己的承認不一致。另一方面，若承認他們准許了割損，則必然會使他陷入另一個異議。因為會有人說：如果宗徒們讚揚你的宣講，卻又准許割損，那他們就是自相矛盾。那麼解決方案是什麼？他能否說他們這樣做是出於對猶太教的遷就？說這話會動搖整個救恩計劃的根基。因此他通過這些話讓這個主題懸而不決：\Quote{但那些被認為是有分量的，與我無關。}這實際上等於說：我不指控也不中傷那些聖人；他們知道他們做了什麼；他們必須向天主交賬。我想要證明的是，他們既沒有推翻也沒有糾正我的做法，也沒有將其視為有缺陷而加以補充，而是給予了認可和同意；弟鐸和巴爾納伯為此作證。然後他補充道：\par

第7節。\BibleQuote{當他們看見我受託傳\OriginalTerm{未受割損者的福音}{Gospel of the Uncircumcision}，正如伯多祿受託傳\OriginalTerm{受割損者的福音}{Gospel of the Circumcision}}——\BibleRef{迦 2:7}\par

\OriginalTerm{受割損者}{the Circumcision}與\OriginalTerm{未受割損者}{the Uncircumcision}，意指的不是這些事物本身，而是由這些區分所標識的民族；因此他補充道：\par

第8節。\BibleQuote{因為那在伯多祿身上運行，使他成為受割損者的宗徒的，也在我身上運行，使我成為外邦人的宗徒。}\BibleRef{迦 2:8}\par

他稱外邦人為未受割損者，稱猶太人為受割損者，並聲明自己的職位與宗徒們相等；並且，通過將自己與他們的領袖而非其他人相比較，他表明兩者的尊嚴是相同的。在他確立了他們意見一致的證據之後，他鼓起勇氣，充滿信心地進行他的論證，不止於宗徒們，而是更進一步，直到基督自己以及祂所賜給他的恩寵，並呼喚宗徒們作他的見證人，說：\par

第9節。\BibleQuote{當他們——雅各伯、刻法和若望，那些被認為是\OriginalTerm{柱石}{pillars}的——看出所賜給我的恩寵時，就向我和巴爾納伯伸出了\OriginalTerm{共融的右手}{right hand of fellowship}。}\BibleRef{迦 2:9}\par

他说不是当他们\Quote{听见}，而是当他们\Quote{觉察}，也就是说，由事实本身得到确证，\Quote{他们就向我和巴尔纳伯伸出了共融的右手}。请注意他是如何逐步证明自己的教训既得到了基督的认可，也得到了宗徒们的认可。因为如果他的宣讲未经基督的认可，恩宠就既不会被植入，也不会在他内运作。当他为了与自己比较而提起时，他单独提到了伯多禄；在这里，当他呼召他们作证时，他一同提到了三人：\Quote{刻法、雅各伯、若望}，并加以赞美：\Quote{那些被视为柱石的人}。这里再次使用\Quote{被视为}的表述并不否定事实的真实性，而是采纳了他人的评价，意味着这些伟大而杰出、名扬四海的人，证明了他的宣讲是经基督认可的，他们实际上是由经验得知并被说服的。\Quote{因此他们将共融的右手}伸向了我，不仅向我，也向巴尔纳伯，\Quote{使我们在外邦人那里去，他们在受割损的人那里去}。这确实展现了极大的智慧，也是他们和好的无可争议的证据。因为这表明，他和他们的教训是可以互换的，两者都赞同同一件事：他们应向犹太人如此宣讲，而他则向外邦人宣讲。因此他补充道：\par

第9节。\BibleQuote{使我们在外邦人那里去，他们在受割损的人那里去。}\BibleRef{迦 2:9}\par

请注意，这里他所说的\Quote{受割损}也不是指割损仪式本身，而是指犹太人；每当他提到这个仪式并希望与之对比时，他会加上\Quote{未受割损}一词。例如他说：\BibleQuote{因为割损固然有益，如果你遵守法律；但如果你违反法律，你的割损就变成了未受割损。}\BibleRef{罗 2:25}又说：\Quote{受割损算不得什么，未受割损也算不得什么。}但当他用这个名称指犹太人而非行为本身，并希望表示民族时，他并非以字面意义的\Quote{未受割损}与之对立，而是以外邦人。因为犹太人与外邦人相对，受割损与未受割损相对。因此他在上面说：\Quote{因为藉着伯多禄在受割损的人中运作的，也藉着我在外邦人中运作}；又说：\Quote{我们在外邦人那里去，他们在受割损的人那里去}，他指的不是仪式本身，而是犹太民族，以此来将他们与外邦人区分。\par

第10节。\BibleQuote{只是他们愿意我们记念穷人；这也是我本来热心去行的。}\BibleRef{迦 2:10}\par

这是他的意思：在我们的宣讲中，我们按天主的旨意将世界分给了我们；我取了外邦人，他们取了犹太人；但对于犹太人中穷人的供养，我也贡献了我的一份；如果我们之间有任何不同意见，他们是不会接受的。其次，这些穷人是谁？许多在巴勒斯坦信主的犹太人被剥夺了全部财产，分散在世界各地，正如他在\WorkTitle{致希伯来人书}中所提到的：\Quote{你们甘心忍受了你们的财物的劫夺}；在写给\WorkTitle{得撒洛尼前书}中，\BibleQuote{你们成了在犹太地为天主而有的各教会的效仿者……因为你们也由你们本族人受了同样的苦，正如他们由犹太人所受的一样。}\BibleRef{得前 2:14}他赞扬了他们的坚韧：并且他在整个过程中表明，那些信主的外邦人并未像信主的犹太人所遭受的那样受到其他人的迫害，因为没有哪个民族如此残忍。因此他表现出极大的热心，正如在\WorkTitle{致罗马人书}\BibleRef{罗 15:25}和\WorkTitle{致格林多人书}\BibleRef{格前 16:1}中显示的那样，这些人应得到许多关注；保禄不仅为穷人募捐，而且亲自送去，如他所说：\BibleQuote{但现在我往耶路撒冷去，供给圣徒}\BibleRef{罗 15:25}，因为他们缺乏生活必需品。在此他表明，这次他决心帮助穷人，已经承担起来，并且不会放弃。\par

通过这些方式表明宗徒们与他之间的同心同德与和谐之后，他不得不继续讲述他在安提约基雅与伯多禄的辩论。\par

第11、12节。\BibleQuote{但当刻法来到安提约基雅时，我当面反对了他，因为他有可责备之处。因为从雅各伯那里来的人还没有到以前，他和外邦人一同吃饭；但他们来到后，他就退避并分离了，害怕那些受割损的人。}\par

许多人因草率阅读书信的这一部分，便以为保禄指责伯多禄虚伪。但事实并非如此，绝非如此，恰恰相反；我们将发现保禄和伯多禄为听众的益处而隐藏在此处的大智慧。但首先必须谈谈伯多禄的直言不讳，以及他如何总是超越其他门徒。确实，正是在这样一次场合中，他因自己信德的刚毅不屈而获得了他的名字。因为当众人都被共同询问时，他抢先一步回答：\BibleQuote{祢是基督，永生天主之子。}\BibleRef{玛 16:16} 当时天国的钥匙交给了他。同样，在山上，他似乎也是唯一的发言者；\BibleRef{玛 17:4} 当基督谈论自己的受难，而其他人保持沉默时，他说：\BibleQuote{这事绝不会临到祢身上。}\BibleRef{玛 16:22} 这些话即使不表明谨慎的性情，至少也表明炽热的爱；在所有实例中，我们发现他比别人更激烈，冲向危险。因此，当基督在岸边被看见，而其他人正推船时，他急不可待，不等船靠岸。\BibleRef{若 21:7} 复活后，当犹太人杀气腾腾、狂怒不止，企图把宗徒们撕碎时，他首先敢于挺身而出，宣告被钉死的那位已被接升天。\BibleRef{宗 2:14} 开启一扇紧闭的门并开始一项行动，比之后畅所欲言更为伟大。一个向如此民众暴露自己生命的人，怎么可能伪装呢？他受鞭打和捆绑时毫不减退勇气，而且在他使命的开始，在充满危险的首都中心——他怎么可能在很久以后，在安提约基雅，那里没有危险，他的品格因行动的见证而光彩熠熠，却害怕那些信主的犹太人呢？我说，他最初在首都并不怕犹太人，反而在许久之后在外地城市，害怕那些已归化的犹太人吗？因此，保禄说这话不是针对伯多禄，而是带着同样的意义，正如他所说：\BibleQuote{那些被认为有些什么的人，无论他们是什么，与我无关。} 但为了消除此点的任何疑虑，我们必须阐明这些说法的理由。\par

宗徒们，如我之前所说，在耶路撒冷准许了割礼，因为突然断绝法律并不可行；但到了安提约基雅，他们不再继续这一惯例，而是与信主的外邦人无分别地生活，此事伯多禄当时也在做。但当有人从耶路撒冷来，他们曾在那里听到他所宣讲的道理，他便不再这样做，担心使他们困惑，而是改变了做法，暗中抱着两个目的：既避免得罪那些犹太人，又给保禄一个合理责备他的借口。因为如果他曾在耶路撒冷传道时允许割礼，而在安提约基雅改变做法，他的行为在那些犹太人看来就像是出于对保禄的恐惧，而他的门徒也会谴责他过分迁就。这会造就不小的绊脚石；但在清楚一切事实的保禄看来，他的退出不会引起这种怀疑，因为他知道其行为的意图。因此，保禄责备，伯多禄顺服，以便当师傅受责却保持沉默时，门徒更容易归顺。没有这一事件，保禄的劝诫效果甚微，但由此提供的严厉申斥的机会，使伯多禄的门徒心中更加恐惧。如果伯多禄反驳保禄的判决，他可能会被正义地指责为破坏安排，但现在一方责备，另一方沉默，犹太派充满了严重惊慌；这就是他为何如此严厉地对待伯多禄。也要注意保禄措辞的细心选择，他以此向有识别力的人指出，他的言辞是出于\OriginalTerm{安排}{\GreekText{οἰκονομίας}}，而非出于愤怒。\par

他的话是：\BibleQuote{当刻法来到安提约基雅时，我当面反对了他，因为他有可责之处；} 也就是说，不是被我，而是被他人所责；如果他本人责难他，他不会不敢这样说。而这话：\BibleQuote{我当面反对了他，} 暗示一种策略，因为如果他们的讨论是真实的，他们就不会在门徒面前彼此责备，因为这对他们会是一个大绊脚石。但现在这个表面的争论对他们大有益处；正如保禄曾向耶路撒冷的宗徒们让步，同样，他们也在安提约基雅向他让步。责备的理由是：\BibleQuote{因为在从雅各伯那里来的几个人来到以前，} 雅各伯是耶路撒冷的导师，\BibleQuote{他常和外邦人一同吃饭，但及至他们来了，他因怕那些受割礼的人，便退避并隔离了自己。} 他恐惧的原因不是自己的危险（因为如果起初他不怕，那时更不会怕），而是他们的背弃。正如保禄自己向迦拉达人所说：\BibleQuote{我为你们害怕，惟恐我在你们身上是白白劳碌了。}\BibleRef{迦 4:11} 又说：\BibleQuote{我只怕你们的心意败坏……，就像蛇用诡诈诱惑了厄娃一样。}\BibleRef{格后 11:3} 因此，他们不知道死亡的恐惧，但害怕门徒丧亡的恐惧却激动他们内心深处。\par

第13节。\BibleQuote{以致连巴尔纳伯也被他们的假装所带走。}\BibleRef{迦 2:13}\par

不要惊奇他把这行为称为\OriginalTerm{掩饰}{dissimulation}，因为正如我之前所说，他不愿透露实情，以便纠正他的门徒。由于他们对法律的热烈依恋，他把当前的行为称为\Quote{\OriginalTerm{掩饰}{dissimulation}，}并严厉责备它，以便有效地根除他们的偏见。\Person{伯多禄}听到后也加入假装，好像他犯了错，以便通过对他施加的责备来纠正他们。如果\Person{保禄}责备这些犹太人，他们会愤怒地拒绝，因为他们轻视他；但如今，当他们看到他们的老师沉默地接受责备时，他们无法鄙视或抗拒\Person{保禄}的判断。\par

第十四节。\BibleQuote{但我看见他们行事不正直，与福音的真理不合。}\BibleRef{迦 2:14}\par

也不要让这句话困扰你，因为他在使用它时并非谴责\Person{伯多禄}，而是为了那些要通过责备\Person{伯多禄}而得到改正的人的利益而如此表达。\par

第十四节。\BibleQuote{我在众人面前对刻法说。}\BibleRef{迦 2:14}\par

请注意他纠正他人的方式：他说\Quote{在众人面前，}以便让听者因此感到警觉。他接着说道——\par

第十四节。\BibleQuote{你既然是犹太人，却按照外邦人的方式生活，而不像犹太人，怎么还勉强外邦人犹太化呢？}\BibleRef{迦 2:14}\par

但和\Person{伯多禄}一同被带走的不是外邦人而是犹太人；那么为什么\Person{保禄}要把未做的事归咎于他，而不是把他的言论指向那些伪装的犹太人，反而指向外邦人？为什么他单独责备\Person{伯多禄}，而其他人也一同伪装？让我们看看他的指控措辞：\Quote{你既然是犹太人，却按照外邦人的方式生活，而不像犹太人，怎么还勉强外邦人犹太化呢？}事实上只有\Person{伯多禄}自己退却了。他的目的是消除对他责备的怀疑；如果他因\Person{伯多禄}遵守法律而责备他，犹太人就会因他对老师的大胆而谴责他。但现在他为他自己的特殊门徒（我指的是外邦人）而控告他，这使他更容易让人接受他所说的话；他也通过避免责备他人，把所有话都指向宗徒而做到这一点。\Quote{你，}他说，\Quote{既然是犹太人，却按照外邦人的方式生活，而不像犹太人；}这几乎等于明确的劝勉要效法他们的老师，他本人虽是犹太人，却按照外邦人的方式生活。然而他并没有这么说，因为他们无法接受这样的建议；但借着责备他以便为外邦人辩护的表象，他揭示了\Person{伯多禄}的真实想法。另一方面，如果他说：你为什么强迫这些犹太人犹太化？他的言辞就会过于严厉。但现在他通过表现出支持外邦门徒而非犹太门徒的一方来实现对他们的纠正；因为适度严厉的责备最容易被人接受。没有一个外邦人能反对\Person{保禄}说他为犹太人辩护。整个困难因\Person{伯多禄}默默接受伪装的指责而消除，以便使犹太人脱离真实情况。起初\Person{保禄}把他的论据指向\Person{伯多禄}所扮演的角色，\Quote{你既然是犹太人：}但他在继续时概括了，并把自己也包括在短语中，\newline{}第十五节。\BibleQuote{我们生来是\OriginalTerm{本来的犹太人}{Jews by nature}，并不是\OriginalTerm{外邦人的罪人}{sinners of the Gentiles}。}\par

第十五节。\BibleQuote{我们生来是\OriginalTerm{本来的犹太人}{Jews by nature}，并不是\OriginalTerm{外邦人的罪人}{sinners of the Gentiles}。}\BibleRef{迦 2:15}\par

这些话是劝勉性的，但由于那些犹太人，被塑造成了责备的形式。同样，在其他地方，他借一个意思传达另一个意思；如他在致罗马人书中所说：\Quote{但我现在往耶路撒冷去，为圣徒服务。}\BibleRef{罗 15:25}这里他的目的不仅是告诉他们他往耶路撒冷去的动机，而是激发他们效法施舍。如果他只想解释动机，只说\Quote{我去为圣徒服务}就够了；但现在请看他还说了什么：\Quote{因为马其顿和阿哈雅乐意凑出一笔捐款，为耶路撒冷圣徒中的穷人。是的，这是他们的乐意，而且他们是欠债者。}又说：\Quote{外邦人既然分享了他们的神恩，也就该在物质上服事他们。}\BibleRef{罗 15:26}\par

请注意他如何抑制犹太人的高傲思想；借一物预备另一物，他的言语带有权威。\Quote{我们生来是\OriginalTerm{本来的犹太人}{Jews by nature}，并不是\OriginalTerm{外邦人的罪人}{sinners of the Gentiles}。}这句话\Quote{本来的犹太人}意味着我们不是皈依者，而是自幼在法律中受教养，我们放弃了习惯的生活方式，转向了在基督内的信仰。\par

第十六节。\BibleQuote{我们知道，人不是因\OriginalTerm{法律的行为}{works of the Law}而成义，而是因对耶稣基督的\OriginalTerm{信德}{faith}；所以我们信了基督耶稣。}\BibleRef{迦 2:16}\par

请在此也注意他多么谨慎地表达；他没有说他们放弃了法律是因为法律是邪恶的，而是因为它软弱。如果法律不能赐予\OriginalTerm{义德}{righteousness}，那么\OriginalTerm{割礼}{circumcision}就是多余的；到目前为止他证明了这一点；但他继续表明它不仅是多余的，而且是危险的。特别值得注意的是，他一开始说人不是因法律的行为而成义；但随着继续，他说话更加强烈；\par

第十七节。\BibleQuote{但如果我们在基督内寻求成义，我们自己也成了罪人，那么基督岂不是\OriginalTerm{罪的仆役}{minister of sin}吗？}\BibleRef{迦 2:17}\par

若信德于祂无益于我们的成义，反而必须再次接受法律，并且若我们因基督的缘故抛弃了法律，却不因此而成义，反因这种抛弃而被定罪，那么我们将发现，那位我们为之抛弃法律而转向信德的基督，竟成了我们定罪的根源。请看，他如何将此事归结为一个必然的荒谬结论。再注意，他辩论时是多么热切而有力。因为他说，若我们不应该抛弃法律，而我们却为基督的缘故抛弃了它，我们便会被审判。那么，你为何向伯多禄强调此事？他比任何人都更了解此事。天主不是曾向他启示，未受割礼的人不应因割礼而被审判吗？他在与犹太人辩论时，不是基于他所见的异象而勇敢地反对吗？他不是从耶路撒冷发出关于此事的明确法令吗？保禄的目的因此并非纠正伯多禄，但责备必须指向他，尽管它针对的是门徒们；不仅针对迦拉达人，也针对其他有同样错误的人。因为虽然现在很少有人受割礼，但通过与犹太人一起守斋和守安息日，他们同样将自己排除于恩宠之外。若基督对仅受割礼的人无效，那么对于守斋和守安息日，从而遵守法律的两条诫命以代替一条的人，就更加危险了。而且时间的因素加重了这一点：起初，当城、圣殿和其他制度尚存时，他们如此行；但这些人，在犹太人受罚、城市毁灭仍历历在目时，比前者遵守更多的法律规条，他们能为这种遵守找到什么借口呢？尤其是在犹太人自己虽极想遵守却无法遵守的时候？你已穿上基督，已成为主的肢体，并被登记在天上的城里，却仍匍匐在法律之下？你怎能获得天国？请听保禄的话：遵守法律摧毁了福音。你若愿意，请学习这事如何发生，并战栗，避开这陷阱。你为何守安息日，并与犹太人一起守斋？是因你惧怕法律及其字句的废弃？但你不会怀有这种恐惧，除非你轻视信德，认为它软弱且独自无力拯救。害怕忽略安息日清楚地表明你仍视法律为有效；若法律是必要的，那么它是整体的必要，而非局部或仅一条诫命；若是整体的，那么因信德而来的义就逐渐被排除。你若守安息日，为何不受割礼？若受割礼，为何不献祭？法律若要遵守，就必须全部遵守，否则就一点也不遵守。若忽略一部分使你惧怕定罪，这惧怕就同样适用于所有部分。若违反整体不受罚，那么违反部分就更不受罚；另一方面，若后者受罚，那么前者更受罚。但若我们必须遵守全部，我们就必须不服从基督，或由于服从祂而成为法律的违犯者。若法律应被遵守，那么不遵守的人就是违犯者，而基督将被发现是这违犯的原因，因为祂自己废除了这些方面的法律，并命令他人废除它。难道你不明白这些犹太化者意欲何为？他们要使那成为我们义之根源的基督成为罪的根源，如保禄所说，\Quote{因此基督成了罪的仆役。}如此将命题归结为荒谬之后，他无需再做任何事来推翻它，仅以此简单的抗议为满足：\par

第17节。\Quote{断乎不可：}因为无耻和不敬无需以推理来对付，仅一个抗议就足够了。 \BibleRef{迦 2:17}\par

第18节。\Quote{若我重建我所拆毁的，我就证明自己是个违犯者。} \BibleRef{迦 2:18}\par

请注意宗徒的辨别力：他的对手试图表明，不遵守法律的人是违犯者，但他将论证回击他们，表明遵守法律的人才是违犯者，不仅是信德的违犯者，而且是法律本身的违犯者。\Quote{我重建我所拆毁的，}即法律；他的意思如下：法律已确然终止，我们已弃绝它，投奔因信德而来的救恩。但若我们坚持要重建它，我们就因这行为本身成为违犯者，试图遵守天主已废除的。接下来他表明法律是如何被废除的。\par

第19节。\Quote{因为我藉着法律已死于法律。} \BibleRef{迦 2:19}\par

这可以从两方面理解：要么他指的是恩宠的法律，因为他常称此为法律，如经上所说：\Quote{因为生命之圣神的法律已使我获得自由；} \BibleRef{罗 8:2} 要么是指旧法律，他说，藉着法律本身，他已死于法律。也就是说，法律本身已教导我不再服从它，因此若我这样做，我将违犯它自身的教导。它如何、以何种方式如此教导？梅瑟论及基督说：\Quote{上主你的天主要从你中间，从你的兄弟中，为你兴起一位像我一样的先知；你们应听从祂。} \BibleRef{申 18:15} 因此，凡不听从祂的人，就是违犯法律。再者，\Quote{我藉着法律已死于法律}这句话也可作另一理解：法律命令其一切规条应被执行，并惩罚违犯者；因此我们都死于它，因为无人完全遵守了它。请注意，他多么谨慎地攻击它；他没有说\Quote{法律已向我死了；}而是说\Quote{我已死于法律；}其意义是：正如一具死尸不可能服从法律的命令，同样，我也如此，因它的诅咒而灭亡，因它的言语而被杀。因此，不要让法律命令死人，死人死于它自己的行为，不仅身体死亡，而且灵魂也死亡，这导致了身体的死亡。他在下文显示了这一点：\par

第19、20节。\Quote{为使我活于天主，我已与基督同钉在十字架上。}\par

说过\Quote{我已经死了}，为避免有人质疑：那么你怎么活着？他补充了活着的原因，并表明活着时法律杀了他，但死后\Person{基督}藉着死亡使他恢复生命。他指出这奇事是双重的：藉着\Person{基督}，那已死的人得以重生，且是藉着死亡。他这里指不朽的生命，因为这是这句\Quote{我活着是为了\Person{天主}，与\Person{基督}同钉十字架}的意思。有人问：一个现在活着呼吸的人如何已被钉十字架？\Person{基督}被钉是明显的，但你如何被钉却仍活着？他这样解释：\par

第20节。\BibleQuote{我还活着，但不再是我，而是\Person{基督}在我内活着。}\BibleRef{迦 2:20}。\par

在这些话中，\Quote{我与\Person{基督}同钉十字架}他暗指圣洗圣事，而在\Quote{然而我活着，但不是我}这句话中，暗指我们随后的生活方式，藉此我们的肢体被钉死。藉着说\Quote{\Person{基督}在我内活着}，他意指没有一件事是我做的，是\Person{基督}不认可的；因为正如藉着死亡，他意指的不是通常理解的死，而是向罪死；同样，藉着生命，他指出从罪中得释放。因为人除非向罪死，否则不能活于\Person{天主}；正如\Person{基督}经历了肉体的死亡，\Person{保禄}也经历了向罪死。\BibleQuote{你们要致死，他说，你们在地上的肢体：淫乱、不洁、邪情}\BibleRef{哥 3:5}，以及，\BibleQuote{我们的旧人已被钉死}\BibleRef{罗 6:6}，这发生在圣洗中。之后，如果你保持向罪死，你就活于\Person{天主}；但如果你让它复活，你就毁坏了你的新生命。然而\Person{保禄}没有这样做，而是保持完全的死；那么，他说，如果我活于\Person{天主}，一种与法律不同的生命，并且向法律死了，我就不可能遵守法律的任何部分。试想他的行走是多么完美，你必会惊叹这位蒙福的灵魂。他并不说\Quote{我活着}，而是说\Quote{\Person{基督}在我内活着}；谁能有胆量说出这样的话？实际上，\Person{保禄}，他已将自己套在\Person{基督}的轭下，并抛弃了一切世俗事物，全面服从祂的旨意，他不说\Quote{我活于\Person{基督}}，而是更为高深的话：\Quote{\Person{基督}在我内活着}。正如当罪掌权时，它本身就是生命力，并将灵魂引向它所愿之处；同样，当罪被治死，而\Person{基督}的旨意被顺从时，这生命就不再是属地的，而是\Person{基督}活着，也就是说，\Person{基督}在我们内工作、掌权。祂说\Quote{我与祂同钉十字架}\Quote{我不再活着}，而是\Quote{死了}，这在许多人看来不可思议，他接着说：\par

第20节。\BibleQuote{并且我现今在肉身内所活的生命，是因信而活，这信德是在\Person{天主}子内的。}\BibleRef{迦 2:20}。\par

他说，上述涉及我们的属灵生命，但这感性的生命，若加以思考，也会发现是因我对\Person{基督}的信德。因为就先前的制度和法律而言，我遭受了最严厉的刑罚，早已灭亡了，\BibleQuote{因为所有人都犯了罪，亏缺了\Person{天主}的荣耀}\BibleRef{罗 3:23}。而我们，本在判决之下，却被\Person{基督}释放了，因为我们所有人都是死的，即便不是实际，至少是按判决；祂将我们从预期的打击中拯救出来。当法律控诉，\Person{天主}定我们的罪时，\Person{基督}来了，藉着舍己至死，将我们从死亡中全部救出。因此，\Quote{我现今在肉身内所活的生命，是因信而活。}若非如此，没有什么能避免一场如洪水般普遍的毁灭，但祂的来临止息了\Person{天主}的义怒，使我们因信而活。这就是祂的意思，从下文可以看出来。在说了\Quote{我现今在肉身内所活的生命，是因信而活}之后，祂接着：\par

第20节。\BibleQuote{在\Person{天主}子内，祂爱了我，并为我舍了自己。}\BibleRef{迦 2:20}。\par

这是怎么回事，\Person{保禄}！你为什么将普遍的恩惠归给自己，将为全世界而作的当作你自己的？因为他不说\BibleQuote{谁爱了我们}，而说\BibleQuote{谁爱了我}。而且福音作者说：\BibleQuote{\Person{天主}如此爱了世界}\BibleRef{若 3:16}；\Person{保禄}自己说：\BibleQuote{祂没有怜惜自己的儿子，但把祂交出}\BibleRef{罗 8:32}——不是只为\Person{保禄}，而是\BibleQuote{为我们众人}；又说：\BibleQuote{为祂自己洁净一个子民，作为祂特有的产业}\BibleRef{铎 2:14}。但考虑到人性的绝望处境，以及\Person{基督}无法言表的温柔关怀——祂从什么中拯救我们，又白白赐给了我们什么——并被对祂的热爱所燃起，他这样表达自己。因此先知们常常将那位万有的\Person{天主}归给自己，如所说：\BibleQuote{\Person{天主}，我的\Person{天主}，我早早寻求你}\BibleRef{咏 62:1}。此外，这语言教导说，每个人都理应向\Person{基督}欠下巨大的感恩债务，好像祂是为了他一个人而来，因为即使只为一人，祂也不会吝啬这屈尊，因此祂对每人的爱之量度与对整个世界一样大。确实，这祭献是为全人类奉献的，足以拯救所有人，但享受这恩惠的只有相信的人。然而，这并没有阻止祂如此伟大的屈尊，因为并非所有人都来；但祂遵循福音中筵席的模式，祂为众人预备了\BibleRef{路 14:16}，然而当客人不来时，祂没有撤去食物，而是召唤了别人。同样，祂也没有轻视那只从九十九只中迷失的羊\BibleRef{玛 18:12}。\Person{圣保禄}也类似地暗示了这一点，当他提到犹太人说：\BibleQuote{即便有些人不信，他们的不信会使\Person{天主}的信实失效吗？断乎不可！不如说\Person{天主}是真实的，人都是虚谎的}\BibleRef{罗 3:3}。当祂如此爱你，以至于舍己，将你——本无希望——带到如此伟大蒙福的生命中，难道你这被恩赐的人，还要回到过去的事物吗？推理完成后，他以强烈的断言结束，说：\par

第21节。\BibleQuote{我不废弛\Person{天主}的恩宠。}\BibleRef{迦 2:21}。\par

让那些今天仍犹太化并固守法律的人听这话，因为这对他们适用。\par

第21节。\BibleQuote{如果义德来自法律，那么基督就白白死了。}\BibleRef{迦 2:21}\par

还有什么罪比这更可憎？还有什么话比这些更令人羞愧？基督的死是法律不能使我们成义的明证；如果法律真能使人成义，那么祂的死便是多余的。然而，这个如此可畏、超越人的理性、不可言喻的奥迹，圣祖们曾为之焦虑，先知们曾预言，天使们曾惊惧地注视，所有人皆承认它是天主温柔的极致——又怎能合理地说它是轻率而徒然完成的呢？保禄思量到，若他们声称如此伟大高深的行为是多余做成的（他们实际的行为正是如此），那将是何等不合情理，因此他甚至对他们使用激烈的言辞，正如接下来的话所显示的。\par

\SourceRecord{Translated by Gross Alexander.；Nicene and Post-Nicene Fathers, First Series；Vol. 13.；Edited by Philip Schaff.；Buffalo, NY: Christian Literature Publishing Co.,；1889.；<http://www.newadvent.org/fathers/23102.htm>.}

% structure: p0001 source=h1 target=section reason=page-title

\section{\WorkTitle{迦拉达书第三篇讲道}}

\BibleRef{迦 3:1} 第一节\par

\Quote{\Emph{愚蠢的迦拉达人！}\Emph{谁迷惑了你们，在你们眼前耶稣基督曾被公开展示、钉在十字架上？}}\par

这里他转入另一个主题；在先前几章中，他已表明自己不是由于人，也不是借着人成为宗徒，更不需要宗徒的教导。如今，既已确立自己作为教师的权威，他便更加自信地论述，并在信德与法律之间作一比较。他在开头曾说：\BibleQuote{我惊奇你们竟这么快就离开了}；\BibleRef{迦 1:6} 但在此处却说：\BibleQuote{愚蠢的迦拉达人}；那时，他的义愤尚在萌发，但如今，在驳斥了对自己的指控并证明之后，这义愤便爆发出来。不要因他称他们为\Quote{愚蠢}而感到惊讶；因为这并非违反基督关于不可称弟兄为愚人的命令，反而正是严格地遵守了这个命令。因为这命令并非简单地说：\BibleQuote{凡向弟兄说\InnerQuote{愚人}的}，\BibleRef{玛 5:22} 而是说，凡如此行的，是\BibleQuote{无故的}。又有谁比他们更适合被称为这样的人呢？他们经历了如此伟大的事件之后，仍固守旧事，仿佛从未发生过别的事。若因此就要称保禄为\Quote{辱骂者}，那么同样，因阿纳尼雅和撒斐辣的事，伯多禄也可被称为杀人者；但若在那件事上如此称呼是荒谬的，那么在此事上更是如此。此外，还应注意，他并非一开始就这样激烈，而是在这些证据和证明之后，如今倒是这些证据和证明本身，而非保禄本人，可以被视为在施行斥责。因为他已表明他们拒绝了信德，并使基督的死亡落空，这时他才引入责备，而这责备即使如此，也比他们所应得的要轻。也请留意他如何迅速收敛：他并没有说：谁诱惑了你们？谁败坏了你们？谁用诡辩迷惑了你们？而是说\BibleQuote{谁嫉妒地看了你们？}这样将责备与某种赞美调和起来。因为这话暗示他们先前的行为曾引起嫉妒，而眼前的事是由于恶魔的恶毒，它的气息毁坏了他们兴旺的境况。\par

当你们在此处听到\Quote{嫉妒}，以及在福音中听到\Quote{恶眼}（两者意思相同）时，切莫认为眼睛的瞥视有天赋能力去伤害被看的人。因为眼睛，即器官本身，不可能是恶的；但基督在那里用这个词指的是嫉妒。简单地观看是眼睛的功能，但以邪恶的方式观看则属于内在堕落的理智。既然通过这一感官，可见之物的知识进入灵魂，而嫉妒多半由财富产生，财富、权势和排场是由眼睛感知的，因此他称眼睛为恶；不是单纯地看，而是出于某种道德堕落而嫉妒地看。因此，借着\BibleQuote{谁曾嫉妒地看你们？}这句话，他暗示这些论者行动并非出于关心，也不是要弥补缺失，而是要毁坏已有的东西。因为嫉妒远非补充所欠缺的，反而减少那些完整的，并败坏了整体。他这样说，并非认为嫉妒本身有能力，而是意指这些教义的传授者出于嫉妒的动机而行。\par

\BibleRef{迦 3:1} 第一节。\BibleQuote{耶稣基督在你们眼前已被公开展示、钉在十字架上。}\par

基督并非在迦拉达，而是在耶路撒冷被钉十字架。他说\BibleQuote{在你们中间}的理由，是要宣示信德看见远处事件的能力。他并非说\BibleQuote{被钉十字架}，而是说\BibleQuote{被公开展示钉在十字架上}，意指借着信德之眼，他们比某些在场亲眼目睹的人看得更清晰。因为后者中许多人并未受益，而前者虽未亲眼看见，却借着信德看得更清楚。这些话既包含赞美，也包含责备：赞美在于他们默默接受了真理；责备在于他们所看见的、为了他们而被剥光、被刺穿、被钉在十字架上、被唾弃、被戏弄、被醋喂、被强盗辱骂、被长枪刺穿的那一位（所有这些都包含在\BibleQuote{被公开展示、钉在十字架上}这句话中），他们竟然离弃了祂，转而去投靠法律，毫不以那些苦难为耻。这里请注意保禄如何略过天地海的一切提述，到处宣讲基督的大能，并像他那样随身带着并高举基督的十字架：因为这就是天主对我们之爱的总和。\par

\BibleRef{迦 3:2} 第二节。\BibleQuote{我只想问你们这件事：你们是因法律的行为，还是因听信德而领受了圣神？}\par

他说：既然你们不留意长篇论述，也不愿深思这救恩计划的伟大，我渴望（鉴于你们极端无知）用简明扼要的论证和总结性的证明来说服你们。此前，他已通过对伯多禄所说的话说服了他们；现在，他完全以出自他们自身事件的论证与他们交锋。他的说服与证据不仅取自别人与你们共享的恩赐，更取自特别赐予你们自己的。因此他说：\BibleQuote{我只想问你们这件事：你们是因法律的行为，还是因听信德而领受了圣神？} 他说：你们已领受了圣神，行了许多大能的事，行了奇迹：复活死人、洁净癞病人、说预言、说方言——法律曾赋予你们这种大能吗？不是信德吗？因为此前你们并不能做这些事。那么，那些从信德领受了如此恩惠的人，竟要离弃它，逃回法律，而法律不能给你们任何同类的东西，这岂不是极端的疯狂吗？\par

\BibleRef{迦 3:3} 第三节。\BibleQuote{你们就这样愚蠢吗？你们既从圣神开始，如今却要由肉体完成吗？}\par

这里他再次适时插入责备；他说，时间本应带来进步；但你们非但没有进步，反而退步了。那些从微小起点开始的人会向更高事物前进；你们始于高处，却跌至低处。即使你们的起点是属血肉的，你们的进步也应是属灵的；但现在，从属灵的事物开始，你们却在属血肉的事物中结束了旅程；因为行奇迹是属灵的，而受割损是属血肉的。奇迹之后你们转向割损，抓住真理后你们退回到预像，凝视太阳后你们寻找蜡烛，有了强健食物后你们奔向奶。他说\Quote{成全了}，这不仅仅意味着\Quote{被引入}，而是\Quote{被祭献}，意指你们的导师像对待动物一样抓住并宰杀你们，而你们却任由自己遭受他们任意摆布。如同某位将军或杰出人物，在千次胜利和战利品之后，竟使自己遭受逃兵的耻辱，并任由他人随意在他身上烙印。\par

4节. \BibleQuote{你们受了这么多的苦，难道徒然了吗？如果真是徒然的话。}\BibleRef{迦 3:4}\par

这句话比前一句更加刺入人心，因为对奇迹的回忆不如他们为了基督而忍受的争战和苦难这样有力。他说，你们所忍受的一切，这些人要使你们失去，还要夺取你们的冠冕。然后，免得他使你们惊慌丧气，他没有直接下判决，而是加上一句：\BibleQuote{如果真是徒然的话}；如果你们只要有心振作精神、恢复自己，他说，那就不是徒然的。那么，那些要切断悔改之路的人在哪里呢？这里有接受了圣神、行了奇迹、成了宣信者、为基督忍受了千般危险和迫害的人，在如此多成就之后从恩宠中坠落了；然而他说，如果你们有此决心，你们可以恢复自己。\par

5节. \BibleQuote{那赐给你们圣神，并在你们中间行奇迹的，是因法律的行为，还是因信德的听闻？}\BibleRef{迦 3:5}\par

他说，你们是否蒙受了如此大的恩赐，成就了如此奇事，是因为遵守法律，还是因为坚持信德？显然是由于信德。既然他们反复争论，认为没有法律信德就无效，他便证明相反的事：即若加上诫命，信德就不再有效；因为信德只有在不加上法律的事物时才有效。\BibleQuote{你们这些靠法律成义的人，已从恩宠上跌落了。}\BibleRef{迦 5:4} 这是他后来在言语更勇敢时说的，利用了那时已获得的有利地位；同时在获得此地位时，他根据他们的经历来论证。因为他说，是在你们服从信德而非法律时，你们才领受了圣神并行了奇迹。\par

这里，既然法律是讨论的主题，他提出了另一个特殊的争论点，而且非常适时且有力地引用了亚巴郎的事迹。\par

6节. \BibleQuote{正如亚巴郎相信了天主，这就算为了他的义德。}\BibleRef{迦 3:6}\par

他说，即使是你们自己所行的奇迹，也宣告了信德的力量；但如果你们允许，我也尝试从古代叙述中引用证据。然后，因为他们非常重视族长，他便举出他的例子，并表明他也是因信德而成义。如果这位在恩宠之前的人，虽拥有许多善工，却因信德而成义，那么我们更是如此。没有法律对他有何损失呢？毫无损失，因为他的信德足以成义。他说，那时法律并不存在，现在也不存在，正如那时一样。在证明不需要法律时，他引用了在法律之前成义的人，免得也有人向他提出反对；因为就如那时法律尚未赐下，同样现在，虽然法律已经赐下，它却被废止了。由于他们非常重视亚巴郎的后裔身份，并且担心放弃法律会被视为与他的血统无关，保禄便通过将他们的论证反过来，消除了这种恐惧，并证明信德尤其与使他们与亚巴郎相连有关。他在《罗马书》中更详细地展开了这一论证；然而他在此也用以下的话进行了敦促：\par

7节. \BibleQuote{所以，你们要知道：具有信德的人，才是亚巴郎的子孙。}\BibleRef{迦 3:7}\par

他通过古代的见证证明这一点：\par

8节. \BibleQuote{圣经预见天主将由信德使外邦人成义，就预先向亚巴郎传报了福音说：万民都要因你获得祝福。}\BibleRef{迦 3:8}\par

这样的话，亚巴郎的子孙就不是那些与他有血缘关系的人，而是那些紧随他信德的人，因为\BibleQuote{万民都要因你}的含义即是如此：显然外邦人被纳入了他的亲属关系。\par

此处也证明了另一重要论点。他们困惑于法律更古老，而信德在后。现在他通过表明信德先于法律来消除这种观念；正如亚巴郎的事例所显明的，他是在法律赐下之前就已成义。他还表明后来发生的事是应验了预言：他说，\BibleQuote{圣经预见天主将由信德使外邦人成义，就预先向亚巴郎传报了福音。}请注意这一点。那位赐予法律的天主，在赐予法律之前，就已决定外邦人要因信德而成义。而且他不说\Quote{启示了}，而说\Quote{传报了福音}，以表示族长对这种成义的方式感到喜悦，并极其渴望它的实现。\par

此外，他们还心存另一种恐惧；经上记载：\BibleQuote{凡不持守法律书上所记载的一切而遵行的，是可诅咒的。}\BibleRef{申27:26}\Person{保禄}极其巧妙而明智地排除了这恐惧，将他们的论点反过来针对他们自己，表明那些放弃法律的人不仅不是受诅咒的，反而是蒙祝福的；而那些遵守法律的人，不仅不是蒙祝福的，反而是受诅咒的。他们说，不遵守法律的人受诅咒，但他证明遵守法律的人受诅咒，而不遵守的人蒙祝福。他们又说，仅持守信德的人受诅咒，但他表明仅持守信德的人蒙祝福。他如何证明这一切呢？因为我们所承诺的并非小事；因此必须仔细留意下文。他先前已经藉着向圣祖所说的话表明了这一点：\BibleQuote{万民都要因你获得祝福。}\BibleRef{创12:4}那时，信德存在，法律尚未存在；所以他总结说：\par

第9节。\BibleQuote{可见，凡是属于信德的人，都同有信德的\Person{亚巴郎}一起蒙受祝福。}\BibleRef{迦3:9}\par

然后，为了不让他们转过来反驳说：不错，\Person{亚巴郎}是因信德而成义，因为那时法律尚未颁布；但法律颁布之后，有什么例子可以证明信德成义呢？\Person{保禄}于是处理这个问题，并证明得比他们所要求的更多：即不仅信德成义，而且法律使它的遵守者处于诅咒之下。要确信这一点，请听宗徒本人的话。\par

第10节。\BibleQuote{凡是出于法律行为的人，都处于诅咒之下。}\BibleRef{迦3:10}\par

这是他在证明之前所作的断言；证明是什么呢？来自法律本身：——\par

第10至11节。\BibleQuote{因为经上记载：\InnerQuote{凡不持守法律书上所记载的一切而遵行的，是可诅咒的。}如今显然，没有人因法律而成义。}\BibleRef{迦3:10-11}\par

因为众人都犯了罪，都处于诅咒之下。但他还尚未如此说，以免显得是出于自己；他再次用一句经文简明扼要地陈述了两点：没有人完全遵守法律（因此他们处于诅咒之下），以及信德成义。那么这句经文是什么呢？是在\Person{哈巴谷}先知书中：\BibleQuote{义人因信德而生活。}\BibleRef{哈2:4}这句经文不仅确立了由信德而来的义德，也表明通过法律无法得救。他说，既然没有人遵守法律，所有人都因过犯而处于诅咒之下，便提供了一条容易的道路，即出自信德的道路；这本身就是无人能因法律成义的有力证明。因为先知不是说：\Quote{义人因法律而生活}，而是：\BibleQuote{因信德}。\par

第12节。\BibleQuote{法律不是出于信德；而是\InnerQuote{遵行它们的人必因此而生活。}}\BibleRef{迦3:12}\par

因为法律不仅要求信德，也要求行为；但恩宠藉着信德拯救并成义。\BibleRef{弗2:8}\par

你看到他是如何证明那些执着于法律的人是处于诅咒之下，因为不可能完全遵守法律；其次，信德如何有这种成义的能力？因为他已经领受的教导，现在以强有力的论据加以坚持。法律太软弱，不能引导人达到义德，于是信德提供了有效的补救，它使得\BibleQuote{法律所不能的事}成为可能。\BibleRef{罗8:3}现在，正如圣经所说：\BibleQuote{义人因信德而生活。}由此驳斥了藉法律得救的观点；更何况\Person{亚巴郎}是因信德而成义，可见信德的效力极其巨大。同样清楚的是：不遵守法律的人受诅咒，而持守信德的人为义人。但也许你问我：我如何证明这诅咒不再有效？\Person{亚巴郎}生活在法律之前，但我们这些曾经屈服于奴役之轭的人，却使自己成为诅咒的罪人；谁能释放我们呢？请看\Person{保禄}对此的迅速回答；他先前的话已经足够；因为一个人一旦成义，并死于法律而拥抱新生命，这样的人如何还能受诅咒呢？然而，这对他来说还不够，于是他开始一个新的论证，如下：——\par

第13节。\BibleQuote{\Person{基督}救赎我们脱离了法律的诅咒，自身成了为我们的诅咒，因为经上记载：\InnerQuote{凡挂在树上的，是可诅咒的。}}\BibleRef{迦3:13}\par

事实上，人们受另一个诅咒的约束，这诅咒说：\BibleQuote{凡不持守法律书上所记载的一切的，是可诅咒的。}\BibleRef{申27:26}我说，人们受这诅咒的约束，因为没有人持守或遵守全部法律；但\Person{基督}用另一个诅咒替换了这个诅咒，即：\BibleQuote{凡挂在树上的，是可诅咒的。}既然挂在树上的人和违反法律的人都受诅咒，而那个将要解除诅咒的人自己必须脱离诅咒，但接受另一个诅咒来代替，因此\Person{基督}承受了这样一个诅咒，从而解除了我们的诅咒。这就像一个无辜的人替一个被判处死刑的人去死，从而把他从惩罚中解救出来。因为\Person{基督}所承受的不是过犯的诅咒，而是另一个诅咒，为的是消除他人的诅咒。因为，\BibleQuote{他没有行过暴力，他口中也没有诡诈。}\BibleRef{依53:9}并且正如藉着死亡，他从死亡中救出了那些正在死去的人，照样，藉着承受诅咒，他使他们脱离了诅咒。\par

第14节。\BibleQuote{为使外邦人也能得到\Person{亚巴郎}的祝福。}\BibleRef{迦3:14}\par

至于外邦人如何呢？经上说：\BibleQuote{在你的后裔中，地上万民都要蒙受祝福：}\BibleRef{创 22:18}这就是说，在基督内。如果这话是对犹太人说的，那么他们本身因犯法而受咒骂，又怎能成为祝福他人的根源呢？一个被咒骂的人，不能将自己所缺少的祝福传递给他人。所以显然，这一切都指向基督，祂是亚巴郎的后裔，外邦人藉祂而得祝福。因此，圣神的恩许也随之而来，正如保禄自己所宣称的：\BibleQuote{使我们能因信德领受圣神的恩许。}因为恩宠不可能降临于无恩宠且犯罪的人身上；他们首先被祝福，咒骂被除去；然后因信德而成义，进而吸引圣神的恩宠。如此，十字架除去了咒骂，信德带来了义德，义德引来了圣神的恩宠。\par

第15节。\BibleQuote{弟兄们，我且按人的常理来说；虽是人的遗嘱，既已确定，也没有人能够废除或加添什么。}\BibleRef{迦 3:15}\par

\BibleQuote{按人的常理来说}是指使用人的例子。他先以圣经、在你们中间行的奇迹、基督的苦难和圣祖为依据，接着转向日常惯例；他向来如此，为要使他的讲论更甘饴，更容易被迟钝之人接受。因此，他与格林多人辩论时：\BibleQuote{谁喂养羊群，而不吃羊群的奶？谁栽种葡萄园，而不吃其中的果子？}\BibleRef{格前 9:7}又如对希伯来人：\BibleQuote{因为遗嘱，必须有立遗嘱者死亡才生效；立遗嘱者活着时，遗嘱从未生效。}\BibleRef{希 9:17}可见他喜爱这类论证。在旧约中，天主也多次这样做，如：\BibleQuote{妇女岂能忘记自己的乳儿？}\BibleRef{依 49:15}又如：\BibleQuote{泥土岂能对塑造它的人说：你做什么？}\BibleRef{依 45:9}在欧瑟亚书中，天主描绘被妻子轻视的丈夫。\BibleRef{欧 2:5}这种使用人的例子也常见于预象中，例如先知取了腰带\BibleRef{耶 13:1}，又下到陶工的家。\BibleRef{耶 18:1}这个例子的意思是，信德比法律更古老，法律是后来的、暂时的，为信德铺路而颁布。因此他说：\BibleQuote{弟兄们，我且按人的常理来说；}前面他称他们为\BibleQuote{愚昧的}，现在称他们为\BibleQuote{弟兄们}，既责备又鼓励。\BibleQuote{虽是人的遗嘱，既已确定。}如果一个人立了遗嘱，他说，谁敢事后推翻或添加什么？这便是\BibleQuote{加添}的意思。那么，天主立约时更是如此；天主与谁立了约呢？\par

第16、17、18节。\BibleQuote{恩许原是向亚巴郎和他的后裔说的。他没有说\InnerQuote{和后裔们}，好像是许多人；而是说\InnerQuote{和你的后裔}，指一个人，就是基督。我是说，天主预先所立定的约，那四百三十年以后才有的法律，不能废除，以使恩许失效。因为如果承继是由于法律，便不再是出于恩许；但天主是藉恩许赐给了亚巴郎。}\BibleRef{迦 3:16}\par

因此，天主与亚巴郎立约，应许在他的后裔中祝福要临到外邦人；这祝福是法律不能废除的。由于这个例子并非完全贴合所论之事，他如此引入：\BibleQuote{我且按人的常理来说}，以免有人从中推导出有损天主威严的事。但让我们深入这个比喻。天主应许亚巴郎，他的后裔要使外邦人蒙福；按肉身来说，他的后裔就是基督；四百三十年后有了法律；如果这法律赐予祝福，即生命和义德，那么应许就失效了。然而，没有人废除人的约，天主的约却在四百三十年后被废除；因为如果赐予应许的不是那约，而是另一个取代它的，那么它就是被废弃了，这是极不合理的。\par

第19节。\BibleQuote{那么法律是什么？它是因为过犯而添上的。}\BibleRef{迦 3:19}\par

这段注释并非多余；请注意，他如何环视一切，仿佛有百眼。他高举了信德，证明了它更古老的权利，但为了防止法律被视作多余，他也纠正了这一方面的教义，证明法律并非无目的地赐下，而是完全有益的。\BibleQuote{因为过犯}，意思是：免得犹太人过于放任，陷入罪恶的深渊，而是将法律放在他们身上作缰绳，引导、规范、阻止他们违犯诫命——至少是部分诫命。因此，法律的裨益不可小觑；但持续多久呢？\par

第19节。\BibleQuote{直到那承受恩许的后裔来到。}\BibleRef{迦 3:19}\par

这是指基督；既然法律给予直到祂来临，你们为何把它延长到自然期限之外呢？\par

第19节。\BibleQuote{并且是藉天使经中保之手设立的。}\BibleRef{迦 3:19}\par

他或是指司祭为天使，或是说明天使本身在颁赐法律时服务。这里的中保他指基督，显示祂在律法之先，并且是律法本身的赐予者。\par

第20节。\BibleQuote{然而中保不是单方面的，但天主是唯一的。}\BibleRef{迦 3:20}\par

对此，异端者能说什么呢？因为按照他们的说法，\Quote{独一真天主}这一表述将圣子排除在真天主之外；同样地，这里的\Quote{天主是唯一}也排除祂身为天主的任何可能。但如果虽然圣父被称为\Quote{唯一天主}，而圣子仍然是天主，那么很明显，虽然圣父被称为\Quote{真天主}，圣子也同样是真天主。现在，他说，中保是在双方之间的；那么基督是谁的中保？显然是天主与人的中保。请注意，他说，基督也赐下了法律；因此，既然法律是祂所赐的，也就属于祂来废除。\par

\BibleRef{迦 3:21} \Quote{那么，法律与天主的恩许相反对吗？}\par

因为如果祝福是在亚巴郎的子孙中赐下的，而法律却带来诅咒，那么它必定与恩许相悖。针对这个异议，他首先以\Quote{绝对不}一词表示抗议，\par

\BibleRef{迦 3:21} \Quote{绝对不：}\par

然后他提出证明：\par

\BibleRef{迦 3:21} \Quote{因为如果曾颁赐一种能赋予生命的法律，那么义德确实出于法律。}\par

他的意思如下：如果我们对生命的望德在于法律，我们的救恩也依赖于它，那么异议或许有效。但如果法律藉着信德拯救你，尽管它使你处于诅咒之下，你却不会因此遭受任何损失，不受伤害，因为信德到来并使一切恢复正确。如果恩许是藉着法律赐予的，你本有理由害怕，离开法律就等于离开义德；但既然恩许是为了禁锢众人，即为了说服众人并揭露他们各自的罪过，那么法律不仅没有将你排除在恩许之外，反而现在帮助你获得恩许。这一点由以下的话表明：\par

\BibleRef{迦 3:22} \Quote{但是圣经把众人都禁锢于罪过之中，为使恩许因对耶稣基督的信德，能赐予相信的人。}\par

由于犹太人甚至对自己所犯的罪过也毫无意识，因而也不渴望赦免；法律被赐下是为了探明他们的伤口，好让他们渴望医生。而\Quote{禁锢}一词意指\Quote{说服}，这种说服使他们陷入恐惧。你看，法律不仅不反对恩许，而且是为恩许而赐下的。如果法律将工作和权柄归于自身，那么异议就成立；但如果它的导向是别的，并且它为此效力，那么它如何与天主的恩许相悖呢？如果法律不曾被赐下，所有人都会因邪恶而丧亡，也就没有犹太人聆听基督了；但如今法律被赐下，它成就了两件事：它使追随者在某种程度上受到德行的训练，并使他们认识到自己的罪过。这尤其使他们更加热切地寻求圣子，因为那些不信的人是由于没有觉察到自己的罪过而不信，正如保禄所表明的：\Quote{因为他们不认识天主的义德，而企图建立自己的义德，他们就没有服从天主的义德。} \BibleRef{罗 10:3}\par

\BibleRef{迦 3:23} \Quote{但在信德来到以前，我们曾被禁锢在法律之下，为那将要启示的信德而被封闭。}\par

在这里，他明确地阐述了我所说的：因为\Quote{我们曾被禁锢}和\Quote{被封闭}这些说法，所指的无非是法律诫命所给予的安全保障；法律像一座堡垒，以恐惧和与之相符的生活将他们围住，从而保守他们归于信德。\par

\BibleRef{迦 3:24} \Quote{因此，法律成了我们的启蒙师，领我们归于基督，好使我们因信德而成义。}\par

原来，\OriginalTerm{启蒙师}{Tutor}并非与\OriginalTerm{导师}{Preceptor}相对，而是与他合作，使青年脱离一切恶习，并有余暇使他为接受导师的教导做好准备。但当青年的习惯养成后，启蒙师便离开他，正如保禄所说。\par

\BibleQuote{但现今信德已经来到，导向完人的境界，我们就不再在启蒙师之下了。因为你们众人藉着对基督耶稣的信德，都是天主的子女。}\par

所以，法律既是我们的启蒙师，我们又曾被禁锢在它之下，它就不是恩宠的敌人，而是恩宠的同工；但如果恩宠已经来到，它却仍将我们束缚，它就变成了敌人；因为如果它阻拦那些本应迈向恩宠的人，那么它就是毁灭我们的救恩。如果一盏在夜间照明的蜡烛，到了白天还阻止我们享受日光，那么它不仅无益，反而有害；同样，如果法律在我们与更大的恩惠之间设置障碍，也是如此。那些仍然持守法律的人正是法律最大的诽谤者，如同启蒙师使一个青年变得可笑，如果时间要求他离开而他却仍将他留在身边。因此保禄说：\BibleQuote{但信德已经来到，我们就不再在启蒙师之下了。}我们不再在启蒙师之下，\BibleQuote{因为你们都是天主的子女。}奇妙！请看信德的能力何其强大，并且他如何随着论述的展开而揭示！先前，他表明信德使他们成为圣祖的子女：\BibleQuote{所以你们要知道，}他说，\BibleQuote{凡属于信德的，就是亚巴郎的子孙；}如今他又证明他们也是天主的子女：\BibleQuote{因为你们众人，}他说，\BibleQuote{藉着信德，就是在基督耶稣内的信德，都是天主的子女；}是藉着信德，而非藉着法律。然后，在说出这件伟大而奇妙的事之后，他还提到了他们被义子的方式：\par

\BibleRef{迦 3:27} \BibleQuote{你们凡受洗归于基督的，就是穿上了基督。}\par

为什么他不说：\Quote{你们凡受洗归于基督的，都是由天主所生的？}因为这才直接证明他们是子女——但他以一种更令人敬畏的方式来陈述：如果基督是天主子，而你们穿上了祂，那么你既把圣子放在你内，又按祂的样式被塑造，你就被带入与祂同一种族和本性。\par

\BibleRef{迦 3:28} \BibleQuote{不再分犹太人或希腊人，不再分为奴的或自由的，不再分男和女，因为你们众人在基督耶稣内已合而为一。}\par

请看何等不知餍足的心灵！因为说了\Quote{我们众人都是藉着\OriginalTerm{信德}{Faith}而成为天主的子女}，他并不就此止步，反倒寻求更精确的表达，以传达与基督更紧密的合一。说了\Quote{你们已穿上了基督}，连这也不够，他更深入这结合，这样注释说：\Quote{你们在基督耶稣内都成了一个}，也就是说，你们都有同一的形像、同一的模型，就是基督的形像。还有什么比这些话更令人敬畏呢？昨天还是希腊人、犹太人，或奴隶的人，如今带着的形像，不是天使或总领天使的，而是万有的主的，并且在自己的身上彰显出基督本人。\par

第29节。\Quote{如果你们属于基督，那么你们就是亚巴郎的后裔，按照恩许是承继人。}\BibleRef{迦 3:29}\par

在这里，你看见，他证实了他先前关于亚巴郎的后裔所陈述的——即恩许是赐给他和他的后裔的。\par

\SourceRecord{Translated by Gross Alexander.；Nicene and Post-Nicene Fathers, First Series；Vol. 13.；Edited by Philip Schaff.；Buffalo, NY: Christian Literature Publishing Co.,；1889.；<http://www.newadvent.org/fathers/23103.htm>.}

% structure: p0001 source=h1 target=section reason=page-title

\section{迦拉达书讲道集 第四篇}

第一节至第三节\BibleRef{迦 4:1-3}\par

\BibleQuote{但是我說，承繼人當還是孩童的時候，他與奴僕毫無分別，雖然他是萬有的主人；但他是在監護人和管家的手下，直到父親所定的日期。我們也是這樣，當我們還是孩童的時候，我們被世俗的小學所奴役。}\BibleRef{迦 4:1-3}\par

這裡的\Quote{孩子}一詞所指的不是年齡，而是理解力；意思是說，天主從起初就為我們預備了這些恩賜，但因我們仍舊\Emph{幼稚}，祂便讓我們受制於世俗的小學，也就是新月和安息日，因為這些日子是根據日月的運行而定的。如果現在他們還強迫你們遵守律法，那就無異於在你們達到完全成熟的時候，又把你們拉回幼稚的狀態。請看遵守節日的後果：主、家主、至高統治者，竟然因此被降到僕人的地位。\par

\BibleQuote{但時期一滿，天主就派遣了祂的兒子，生於女人，生於法律之下，為要把在法律之下的人贖出來，使我們獲得義子的地位。}\BibleRef{迦 4:4-5}\par

這裡他提到了降生成人的兩個目的和效果：脫離邪惡和賜予美善，這些是唯獨基督才能成就的。那就是：脫離法律的詛咒，並提升到兒子的身份。他恰當地說，我們可以\Quote{獲得}\Quote{領受}，表示這是應得的；因為天主從前曾向亞巴郎應許了這些事，正如宗徒自己在別處詳盡論述的。我們怎樣成為兒子的呢？他告訴了一個方式：我們穿上了那為兒子的基督；現在他又提到另一個：我們領受了義子的聖神。\par

\BibleQuote{因為你們是兒子，天主就派遣了祂兒子的聖神進入你們心中，呼喊：\InnerQuote{\OriginalTerm{阿爸}{Abba}，父啊！}所以你們不再是奴僕，而是兒子；如果是兒子，也藉著天主成為承繼人。}\BibleRef{迦 4:6-7}\par

如果我們沒有先成為兒子，我們就不能稱祂為父。那麼，如果恩寵已經使我們從奴僕變成自由人，從孩童變成成人，從外人變成承繼人和兒子，我們卻離棄這恩寵，向後倒退，豈不是極其荒謬和愚蠢嗎？\par

\BibleQuote{但那時你們不認識天主，就給那些本來不是神的作了奴僕。現在你們認識了天主，更可說是被天主所認識，怎麼還要再轉回那軟弱無用的小學，情願再作它的奴僕？}\BibleRef{迦 4:8-9}\par

這裡他轉向信主的外邦人，說這種嚴守節日是一種偶像崇拜，如今要受重罰。為了強調這點並激發他們更深的警覺，他稱那些要素為\Quote{本來不是神}。他的意思是：那時你們因在黑暗中迷惘，趴在地上；但如今你們認識了天主，更可說是被天主所認識，如果經過這樣的對待後，你們又陷入同樣的毛病，你們將為自己招來何等重大嚴厲的懲罰！你們不是靠自己的努力找到天主，而是在你們還在錯謬中時，祂吸引你們到祂自己那裡。他說\Quote{軟弱無用的小學}，因為它們對我們所期待的福樂毫無助益。\par

\BibleQuote{你們謹守日子、月份、節期、年份。}\BibleRef{迦 4:10}\par

由此可見，他們的教師不僅向他們宣講割損，也宣講節日和月朔。\par

\BibleQuote{我為你們害怕，唯恐我對你們的勞苦歸於徒然。}\BibleRef{迦 4:11}\par

請看宗徒的溫情關懷；他們動搖了，他就顫抖恐懼。因此他這樣說以徹底使他們羞愧：\Quote{我對你們的勞苦}，意思是，不要使那曾令我流汗痛苦的努力化為虛空。他說\Quote{我害怕}，並加上\Quote{唯恐}一詞，既激起警覺，又鼓勵好的希望。他不是說\Quote{我已經徒然勞苦}，而是說\Quote{唯恐}，這等於說：沉船尚未發生，但我看見大風暴將至；因此我害怕，但並非絕望；你們有能力扭轉一切，回到從前的平靜。接著，他彷彿向這些在風暴中的人伸出手來，將自己置身於他們中間，說：\par

\BibleQuote{弟兄們，我請求你們，要像我一樣，因為我也像你們一樣。}\BibleRef{迦 4:12}\par

這是對他猶太門徒說的，他舉自己的榜樣來勸導他們放棄舊習俗。他說，即使你們沒有其他榜樣，只要看看我就足夠使你們改變並獲得勇氣。因此，注視我；我也曾一度和你們一樣心態，尤其如此；我對法律有火熱的熱忱；但後來我毫不猶豫地放棄了法律，脫離了那種生活方式。你們非常清楚我是如何頑固地緊抓住猶太教，又是如何更強烈地將它放下。他將這一點放在最後是恰當的：因為大多數人，即使給他們上千個理由，而且都是正當的，卻更容易被與自己相似的情況影響，且更堅定地持守他們看到別人做過的事。\par

\BibleQuote{你們並沒有虧負我。}\BibleRef{迦 4:12}\par

請看他如何再次用榮譽的稱呼稱呼他們，這也提醒他們恩寵的道理。他嚴厲地責備了他們，從各方面收集證據，指出他們違反法律，多方面打擊他們之後，他退讓並安撫他們，說話更加溫和。因為如果只會安撫就會導致疏忽，而如果總是尖銳地責備會使人煩躁；所以必須處處保持適當的比例。請看他如何為自己所說的話辯解，表明那並非出於不喜歡他們，而是出於關切。在狠狠地割了一刀之後，他像倒油一樣傾注這種鼓勵；並且表明他的話不是出於仇恨或敵意，他提醒他們對他所表現的愛，將自辯與讚美混合在一起。因此他說：\Quote{你們並沒有虧負我}。\par

\BibleQuote{你們知道我初次傳福音給你們，是因為身體有疾病。你們並沒有輕視或厭棄我在肉體上的考驗。}\BibleRef{迦 4:13-14}\par

不伤害人确实不算什么大事，因为没有人会无缘无故地、毫无目的地去烦扰一个从未伤害过他的人。但至于你们，你们不仅没有伤害我，反而向我显示了极大且难以言喻的善意，一个受到如此礼遇的人不可能出于恶意说出这样的话。那么我的言语不可能源于恶意；由此看来，它出自关爱与挂念。\Quote{你们没有亏负我；你们知道，我因肉身的软弱曾向你们宣讲福音。}还有什么比这神圣的灵魂更温柔、更甘甜、更充满情谊的呢！他之前所说的那些话，并非出于无理的愤怒，也不是激情的冲动，而是出于极大的挂虑。我为什么说你们没有伤害我？你们反而向我表现出了巨大而真诚的关怀。因为\Quote{你们知道}，他说，\Quote{我因肉身的软弱曾向你们宣讲福音；并且你们对我肉身的考验，没有轻视，也没有拒绝。}他是什么意思？当我向你们宣讲时，我被人驱赶、被人鞭打、历经千辛万苦，而你们并没有鄙视我；这就是\Quote{你们对我肉身的考验，没有轻视，也没有拒绝}的意思。请看他的属灵智慧；在自我辩护中，他再次通过展示自己为他们所受的苦来打动他们的情感。他说，这事并没有冒犯你们，你们也没有因我的苦难和迫害而拒绝我；或者如他现在所称的，他的软弱和考验。\par

第十四节：\BibleQuote{你们却接待我，如同天主的使者。}\BibleRef{迦 4:14}\par

那么，当他们迫害、驱赶他时，却接待他如同天主的使者；而当他敦促他们行合适之事时，反而不接待他，这岂非荒谬？\par

第十五、十六节：\Quote{那么，你们当初的庆幸在哪里呢？我给你们作证，如果可能，你们甚至会把眼睛剜出来给我。难道因为我向你们说实话，就成了你们的仇敌吗？}\par

此处他显出困惑与惊奇，并希望从他们自己那里得知态度转变的原因。他说：谁欺骗了你们，使你们对我的态度发生了改变？你们不是曾经服事我，视我比自己的眼睛更宝贵吗？发生了什么？这厌恶从何而来？这猜疑从何而起？是因为我向你们说了实话吗？正因如此，你们更应加倍尊敬我、关注我；然而，\Quote{我因向你们说实话，反而成了你们的仇敌}——因为我再也找不到其他原因。请看他自我辩护中是何等谦虚；他没有用自己的行为向证明，而是用他们对他的行为来证明，他的言语不可能源于恶意。因为他不是说：一个为你们受鞭打、被驱赶、经受过无数苦难的人，怎么会如今图谋害你们呢？而是从他们引以为豪的事出发，说：一个受你们尊敬、被你们当作天主的使者接待的人，怎么会以相反的行为来回报你们呢？\par

第十七节：\Quote{他们热心追求你们，并非出于善意；他们反而想把你们关在外面，好让你们去追求他们。}\BibleRef{迦 4:17}\par

一种有益的竞赛引导人效法德行，但一种恶意的竞赛却引诱行在正路上的人偏离德行。这正是那些人的目的，他们想剥夺你们圆满的知识，将残缺的、虚假的知识灌输给你们，目的无他，就是要占据教师的位置，并把你们这些本来比他们更高的人贬为门徒。因为\Quote{好让你们去追求他们}这句话就是这个意思。但他说，我渴望相反的情况，愿你们成为他们的榜样，成为更高完美的典范；这事在我与你们同在时确实发生过。因此他接着说，\par

第十八节：\Quote{但在好事上被人热心追求，常常是好的，不仅是我与你们同在的时候。}\BibleRef{迦 4:18}\par

此处他暗示，他的缺席是此事的原因，真正的福分是门徒不仅在师傅面前，也在师傅不在时持守正确的见解。但因他们尚未达到这完美境界，他便竭尽全力帮助他们达到。\par

第十九节：\Quote{我的孩子们，我为你们再受产痛，直到基督在你们内成形。}\BibleRef{迦 4:19}\par

请看他的困惑与不安：\Quote{弟兄们，我恳求你们}；\Quote{我的孩子们，我为你们再受产痛}。他就像一位为子女颤抖的母亲。\Quote{直到基督在你们内成形}。请看他的慈父柔情，请看这配得宗徒的沮丧。听他发出的哀叹，比分娩的妇人更加尖锐——你们毁坏了肖像，破坏了亲情，改变了形态，需要再一次的再生和重塑；尽管如此，我仍称你们为儿女，即便你们是畸形的怪物。但他没有这样说，而是宽恕他们，不愿打击他们，不愿在伤口上加伤。明智的医生不会一下子治好长期患病的人，而是逐渐治疗，以免病人昏厥而死。这位有福的人也是如此；因为他心中的痛苦随爱之深切而更加剧烈。这过错并非小事。正如我一直说并将永远说的，即便轻微的瑕疵也会破坏整体的外貌，扭曲整个形象。\par

第二十节：\Quote{我恨不得现在就与你们同在，改变我的声调。}\BibleRef{迦 4:20}\par

请看他的热情，他无法自制，无法隐藏这些情感；这就是爱的本性。他不满足于言词，而是渴望与他们同在，以便如他所说，改变他的声调，即转为哀哭，流泪，将一切都变为哀悼。因为他无法通过书信表达他的眼泪或悲痛的呼喊，因此他热切地渴望与他们同在。\par

20节。\BibleQuote{我为你们而感到困惑}\BibleRef{迦 4:20} 我不知道该说什么、该想什么。你们这些曾为信德冒险、并通过信德行奇事、升到最高天的人，怎么会突然堕落到如此深的地步，被引向割损或安息日，完全依赖犹太派呢？因此，他在开头就说：\BibleQuote{我惊讶你们这么快就离开了}，而这里他说：\BibleQuote{我为你们而感到困惑}，好像说：我该说什么？我该讲什么？我该想什么？我深深困惑。所以他必然要哭泣，如同先知们在困惑时一样；因为不仅劝诫，而且哀哭也常是关怀的表现。他在\Place{米肋托}的讲话中说：\BibleQuote{三年之久，我日夜不停地流泪劝诫你们每一个人}\BibleRef{宗 20:31}，在此也说：\BibleQuote{并改变我的声音}。当我们在出乎意料的困惑和无助中时，我们被催逼流泪；所以\Person{保禄}严厉地劝诫他们，并试图使他们羞愧，然后又安慰他们，最后他哭泣。这哭泣不仅是责备，也是安抚；它不像责备那样激怒人，也不像纵容那样放松，而是一种混合的疗法，在劝诫中非常有效。这样，他通过眼泪软化了他们的心，有力地抓住了他们的心，然后再次进入争战，提出了一个更大的命题，证明法律本身也反对遵守它。之前，他举了\Person{亚巴郎}的例子，但现在（更有力的是）他引用了法律本身，命令他们不要遵守它，而是放弃它。所以，他说：你必须放弃法律，如果你要服从它，因为这是它自己的意愿；但他没有明说，而是以另一种方式强调，混合了事实的叙述。\par

21节。他说：\BibleQuote{请告诉我，你们这些愿意在法律之下的人，难道没有听见法律吗？}\BibleRef{迦 4:21}\par

他说得对：\Quote{你们这些愿意（在法律之下）的人}，因为这件事不是出于正常有序的传承，而是他们自己不合时宜的好争论。他这里称\WorkTitle{创世纪}为法律，他常以此名称称呼整个旧约。\par

22节。\BibleQuote{因为经上记载：\BibleRef{创 15:16} \Person{亚巴郎}有两个儿子，一个由婢女所生，一个由自由妇人所生。}\BibleRef{迦 4:22}\par

他再次回到\Person{亚巴郎}，不是重复，而是因为这位圣祖在犹太人中的声望很大，为要表明\OriginalTerm{预像}{type}源自那里，现在的事件早已在他身上预表。他先前已表明迦拉达人是\Person{亚巴郎}的儿子，现在，由于这位圣祖的儿子们地位不等，一个由婢女所生，一个由自由妇人所生，他表明他们不仅是他的儿子，而且像那生来自由高贵的人一样是他的儿子。信德的力量就是如此。\par

23节。\BibleQuote{但那婢女的儿子是按血肉生的，那自由妇人的儿子却是藉恩许生的。}\BibleRef{迦 4:23}\par

\Quote{按血肉}是什么意思？他说信德使我们与\Person{亚巴郎}结合，而听众觉得不可思议：那些不是\Person{亚巴郎}所生的人怎么被称为他的儿子？他证明这个悖论早已发生；因为\Person{依撒格}不是按自然秩序、婚姻法则或血肉能力所生，却真是他的儿子。他是出自已死的身体和已死的子宫；他的成胎不是由于血肉，他的出生不是由于种子，因为子宫因年龄和不育而死亡，但天主的圣言塑造了他。那婢女之子则不然；他是按自然法则、依婚姻方式而生的。然而，那不是按血肉所生的比按血肉所生的更尊贵。所以，不要因你们不是按血肉所生而困扰；正因你们不是这样生的，你们才最是\Person{亚巴郎}的亲属。按血肉所生并不使人更尊贵，反而更不尊贵，因为不是按血肉所生的出生更奇妙、更属神。从古时生的人可以明显看出；例如，\Person{依市玛耳}是按血肉所生的，不仅是奴仆，还被逐出父家；但\Person{依撒格}是按恩许所生的，是真儿子、自由人，是万物的主人。\par

24节。\BibleQuote{这些都是\OriginalTerm{寓意}{allegory}。}\BibleRef{迦 4:24}\par

与通常用法相反，他称\OriginalTerm{预像}{type}为\OriginalTerm{寓意}{allegory}；他的意思如下：这段历史不仅宣告表面上的事，还宣布了更进一步的事，因此被称为\OriginalTerm{寓意}{allegory}。它宣布了什么？正是现在所有的事。\par

24节。他说：\BibleQuote{这两个妇人代表两个盟约；一个是出自\Place{西乃山}，生子为奴，那就是\Person{哈加尔}。}\BibleRef{迦 4:24}\par

\Quote{这些：}指谁？那些孩子的母亲，\Person{撒辣}和\Person{哈加尔}；它们是什么？两个盟约，两个法律。正如历史中给了妇人的名字，他沿用这两个种族的称号，表明从名字本身可以推出多少事。如何从名字推出？\par

25节。\BibleQuote{这\Person{哈加尔}是指\Place{阿拉伯}的\Place{西乃山}。}\BibleRef{迦 4:25}\par

婢女名叫\Person{哈加尔}，\Quote{哈加尔}在那国语言中是\Place{西乃山}的词。所以所有出生于旧盟约的人必定是奴仆，因为颁布旧盟约的那座山与婢女同名。它还包括\Place{耶路撒冷}，因为这就是下面句子的意思：\par

25节。\BibleQuote{与现在的\Place{耶路撒冷}相对。}\BibleRef{迦 4:25}\par

也就是说，它毗邻并接连于它。\par

25节。\BibleQuote{因为她与她的儿女同为奴仆。}\BibleRef{迦 4:25}\par

由此得出什么结论？不仅她为奴并生了奴仆，而且这个盟约也是如此，婢女是它的\OriginalTerm{预像}{type}。因为\Place{耶路撒冷}毗邻与婢女同名的山，而在这座山上颁布了盟约。那么，\Person{撒辣}的\OriginalTerm{预像}{type}在哪里？\par

26节。\BibleQuote{但那在上\Place{耶路撒冷}是自由的。}\BibleRef{迦 4:26}\par

因此，凡由她而生的人不是奴仆。这样，下方耶路撒冷的预像就是哈加尔，正如那山被称为哈加尔所表明的；但上面的耶路撒冷的预像是教会。然而，他不满足于这些预像，还加上依撒意亚的见证来支持他所说的话。他曾说上面的耶路撒冷\Quote{是我们的母亲}，并将这名号给了教会，然后他引用了先知的支持：\par

第27节。\BibleQuote{荒胎的，不生育的啊，欢跃吧！未曾经过产痛的啊，欢呼高唱吧！因为被遗弃者的子女比有夫之妇的子女更多。}\BibleRef{迦 4:27}\BibleRef{依 54:1}\par

这先前被称为\Quote{荒胎}和\Quote{孤寂}的是谁？显然，它就是外邦人的教会，那从前被剥夺了对天主的认识的吗？而那\Quote{有丈夫的}是谁？显然是会堂。然而，荒胎的妇女在子女数目上超过了有夫的，因为后者只包含一个民族，但教会的儿女却充满了希腊人和外邦人的国土、大地和海洋，整个可居住的世界。请注意，撒辣如何通过行为，先知如何通过言语，描述了将要发生在我们身上的事。也请注意，依撒意亚称为荒胎的那一位，保禄证明她有许多子女，这在撒辣身上也已经预像性地发生了。因为她尽管荒胎，却成了众多后代之母。然而保禄不满足于此，他仔细追溯荒胎妇女成为母亲的方式，以便在这细节上预像也与真理相符。因此他接着说道：\par

第28节。\BibleQuote{弟兄们，我们像依撒格一样，是恩许的子女。}\BibleRef{迦 4:28}\par

不仅教会像撒辣一样荒胎，或像她一样成为多子之母，而且她生子的方式也与撒辣相同。因为使撒辣成为母亲的不是自然，而是天主的恩许，[因为天主的话语曾说过：\BibleQuote{到约定的时期，我必回到你这里，撒辣将有一个儿子。}\BibleRef{创 18:14}这话进入母腹形成了婴孩；]同样，在我们的重生中，使人成形的不是自然，而是由司铎宣讲的天主之言（信友知道它们），这言在水洗中，如同在一个母腹内，形成并重生受洗的人。\par

因此，如果我们都是荒胎妇人的儿子，那么我们是自由的。但有人可能反驳说：这是一种什么样的自由，当犹太人抓住并鞭打信徒，那些拥有这自由表象的人却受迫害？因为这些事情当时在迫害信友时发生了。他回答说：不要让这事困扰你们，这在预像中也已预见，因为自由的依撒格曾受到为奴的依市玛耳迫害。因此他接着说：\par

第29至30节。\BibleQuote{但是当时那按肉身而生的迫害那按圣神而生的，如今也是这样。然而圣经怎么说？\BibleRef{创 21:10}把婢女和她的儿子赶出去，因为婢女的儿子不能与自由妇人的儿子一同承受产业。}\BibleRef{迦 4:29-30}\par

什么！这全部的安慰就在于表明自由人被奴仆迫害吗？绝不然，他说，我不就此止步；听接下来所说的，然后，如果你们在迫害之下不胆怯，你们将得到充分的安慰。接下来是什么呢？\BibleQuote{把婢女的儿子赶出去，因为他不能与自由妇人的儿子一同承受产业。}\BibleRef{创 21:10}看，一时暴虐和不合时宜的鲁莽的结果！那儿子被他父亲的房屋赶出，与他的母亲一起成了被放逐的漂泊者。也请思考这话的智慧；因为他不是说仅仅因为迫害而被赶出，而是说他不能成为继承人。因为这惩罚不是因着他暂时的迫害而加于他（因为那本是小事，无关宏旨），而是他不被允许参与为儿子预备的产业。这证明，撇开迫害不谈，这件事从一开始就被预像了，并非起源于迫害，而是出自天主的旨意。他也没有说：\Quote{亚巴郎的儿子不能成为继承人}，而是说：\Quote{婢女的儿子}，以较低微的出身来区分他。现在，撒辣是荒胎的，外邦人的教会也是如此；请注意预像如何在其每一细节上得以保存：如前者在过去的年岁中一直不怀孕，并在极其老迈时成为母亲，同样，后者在时期一满时便生产。而先知们早已宣告了这事，说：\BibleQuote{荒胎的，不生育的啊，欢跃吧！未曾经过产痛的啊，欢呼高唱吧！因为被遗弃者的子女比有夫之妇的子女更多。}\BibleRef{依 54:1} 他们在此指的是教会；因为她不认识天主，但一旦认识了他，便超过了多产的会堂。\par

第31节。\BibleQuote{弟兄们，我们不是婢女的子女，而是自由妇人的子女。}\BibleRef{迦 4:31}\par

他转身从各个方面讨论这事，想要证明所发生的事并非新奇，而是在许多世代以前已被预像了。那么，那些被分别出来如此之久、并获得自由的人，自愿服从奴役的轭，岂非荒谬至极？\par

接着，他提出了另一个劝诱，要他们坚持他的教训。\par

\SourceRecord{Translated by Gross Alexander.；Nicene and Post-Nicene Fathers, First Series；Vol. 13.；Edited by Philip Schaff.；Buffalo, NY: Christian Literature Publishing Co.,；1889.；<http://www.newadvent.org/fathers/23104.htm>.}

% structure: p0001 source=h1 target=section reason=page-title

\section{《迦拉达书》第五篇讲道}

第一节 \BibleRef{迦 5:1}\par

\Quote{\Emph{基督以自由解放了我们；所以要站稳。}}\par

难道你自己完成了释放，竟要再回到从前所受的统治之下吗？是另一位救赎了你，是另一位为你付了赎价。请注意，他用了多少方法引导他们远离犹太主义的错误：首先表明，那些由奴隶变为自由的人，竟渴望再成为奴隶而非自由，这是极端的愚蠢；其次，他们将因忽视和忘恩负义对待恩主而被定罪，因他们藐视那解救他们的，却爱那奴役他们的；第三，这是不可能的。因为那一位一次为我们全部赎了身，法律就不再有任何权柄。通过\Quote{站稳}这个词，他指出他们的动摇。\par

第一节 \BibleRef{迦 5:1} \Quote{不要再被奴役的轭缠住。}\par

通过\Quote{轭}这个词，他指出这种行为的沉重负担；通过\Quote{再}这个词，他指出他们的完全无知。如果你们从未体验过这负担，就不会受到如此严厉的责备；但你们通过尝试已经知道这轭多么令人厌烦，却还要再屈服于它，这实在是不可宽恕的。\par

第二节 \BibleRef{迦 5:2} \Quote{看，我保禄告诉你们：如果你们接受割损，基督就对你们毫无益处。}\par

看，多么大的威胁！所以他合理地诅咒了天使。那么基督如何对他们毫无益处呢？因为他没有用论证来支持这一点，只是宣告了它，他的权威所应得的信任，仿佛弥补了所有后续的证明。因此他开篇就说：\Quote{看，我保禄告诉你们}，这是一个对他所断言充满信心的人的表达。我们将尽可能附加我们自己的理由，说明基督对那些受割损的人如何毫无益处。\par

受割损的人是出于对法律的恐惧而受割损；畏惧法律的人，就不信任恩宠的能力；不信任的人，就不能从所不信任的事物中获得任何益处。或者这样看：受割损的人使法律生效；但既然认为它有效，又在大部分上违反它，只在较小部分上遵守它，他就再次使自己处于诅咒之下。一个使自己服从诅咒，又拒绝因信德而来的自由的人，怎能得救呢？如果可以说出看似矛盾的话，这样的人既不信基督，也不信法律，而是站在两者之间，想要同时从两者获益，结果却从两者都得不到果实。在说了基督对他们毫无益处之后，他简短而精辟地提出了证明，如下：\par

第三节 \BibleRef{迦 5:3} \Quote{我再向每一个受割损的人作证：他有义务遵守全部法律。}\par

为使你们不认为这是出于恶意，他说，我不仅对你们说，而且对每一个受割损的人说：他有义务遵守全部法律。法律的各部分相互联结。正如一个从自由人自愿成为奴隶的人，不再随心所欲，而是受所有奴隶法律的约束；同样，在法律方面，如果你承担它的一小部分，并屈服于这轭，你就引来了它的全部统治。世俗的继承也是如此：一个不碰遗产的人，就免除了由继承死者而来的所有事务；但如果他拿了一小部分，即使不是全部，也因那部分而须对一切负责。这情况不仅在我提到的方式中发生，而且也在另一种方式中发生，因为法律的规条是相互关联的。例如：割损与祭献相连，也与遵守日子相连；祭献又与遵守日子和地点相连；地点又涉及无尽的洁净细节；洁净又带来一大串各种各样的规条。因为不洁的人不可以献祭，不可进入圣所，也不可行任何其他类似的事。因此，法律甚至通过一条诫命就引入了许多东西。如果你们受割损，但不是在第八天，或者在第八天，但没有献祭，或者献了祭，但地点不合规定，或者地点合规定，但对象不合常规，或者对象合常规，但你们不洁，或者洁净却没有按规矩净化，那么一切都归于无效。因此他说\Quote{他有义务遵守全部法律}。如果法律有效，就不要只履行一部分，而要履行全部；但如果法律无效，就连一部分也不可履行。\par

第四节 \BibleRef{迦 5:4} \Quote{你们这些要靠法律成义的人，是与基督隔绝了；从恩宠中跌落了。}\par

在确立了他的论点之后，他终于宣布他们面临最严厉惩罚的危险。当一个人诉诸于不能拯救他的法律，并从恩宠中跌落后，除了不可挽回的报应之外，还剩下什么？法律无能为力，恩宠又拒绝了他。\par

这样，他加剧了他们的恐惧，扰乱了他们的心，向他们展示了他们即将遭遇的全部沉船，又为他们打开了近在咫尺的恩宠港湾。这是他的一贯做法，他表明在此方面救恩是容易而稳妥的，接着说道：\par

第五节 \BibleRef{迦 5:5} \Quote{因为我们藉着圣神，因信德，等待正义的希望。}\par

他说，我们不需要那些法律的规条；信德足以使我们获得圣神，并且藉着祂获得正义，以及许多伟大的恩惠。\par

第六节 \BibleRef{迦 5:6} \Quote{因为在耶稣基督内，割损或不割损都毫无功效；惟有以爱德行事的信德，才算什么。}\par

注意他如今以何等大的勇气面对他们。他说：凡穿上基督的人，不要再为这等事操心。他先前说割损是有害的，如今又何以视之为无关紧要？对于在信德之前已有割损的人，它是无关紧要的；但对于在信德赐下之后才受割损的人，则并非如此。也要注意他是如何将其与未受割损并列，借此阐明它的观点；是信德造成了区别。如同在选择角斗士时，他们是钩鼻还是扁鼻，黑或白，在比赛中都没有影响，只需寻找强壮而技巧娴熟的人；同样，所有这些身体上的偶然，对那将列入新盟约的人既不造成损害，也无帮助。\par

\Quote{藉着爱德行事}是什么意思？这里他给了他们沉重的打击，指出这错误的侵入是由于基督的爱未在他们心中扎根。因为不仅需要相信，还需持守于爱。他仿佛说：如果你们爱基督当如所爱，就不会离弃奴役，也不会抛弃那救赎你们的主，更不会侮慢那赐给你们自由者。这里他也暗示那些图谋反对他们的人，意指他们若对这些人怀有友爱，就不敢如此行。他也愿藉这些话纠正他们的生活态度。\par

第7节：\BibleQuote{你们跑得好，是谁阻拦了你们？}\BibleRef{迦 5:7}\par

这不是质问，而是困惑与悲伤的表达。这样的进程如何被中断？谁有如此能力？你们超越一切，居于教师之列，却连门徒的位置都无法保持。发生了什么？谁能做到？这更像是一个呼喊与哀恸之人的话，正如他先前所说的：\Quote{谁迷惑了你们？}\BibleRef{迦 3:1}\par

第8节：\BibleQuote{这劝诱不是出自那召你们者。}\BibleRef{迦 5:8}\par

那召你们者，召你们并非为了这样的变动，他没有定下法律，要你们犹太化。然后，恐无人反驳：\Quote{你为何如此夸大其词，加重事态？我们只守了律法的一条诫命，你却如此大声喧哗？}请听他如何以未来而非现今之事恐吓他们，话语如下：\par

第9节：\BibleQuote{一点面酵能使整个面团发酵。}\BibleRef{迦 5:9}\par

他说，这轻微的错误若不纠正，就有能力（如酵与面团）将你们引向完全的犹太主义。\par

第10节：\BibleQuote{我在主内对你们深信，你们必不会心存别样。}\BibleRef{迦 5:10}\par

他并没有说：\Quote{你们不是这样想的}，而是说：\Quote{你们必不会这样想}；也就是说，你们将被纠正。他如何知道？他并非说：\Quote{我知道}，而是说：\Quote{我信赖天主，并祈求祂的助佑以纠正你们，我怀有希望}；他不仅说：\Quote{我对主有信心}，而且说：\Quote{我在主内对你们有信心}。他处处将责备与赞美相连；这里仿佛他说：我知道我的门徒，我知道你们乐于被纠正。我抱有良好的希望，一方面是由于主，他不允许任何事物（即使微不足道）丧失，另一方面是由于你们，你们能迅速复原。同时他劝勉他们自己努力，因为如果我们不贡献自己的努力，就不可能从天主获得帮助。\par

第10节：\BibleQuote{但那扰乱你们的人，无论他是谁，必承担他的审判。}\BibleRef{迦 5:10}\par

他不仅用鼓励的话语，而且通过诅咒或预言反对他们的教师，以此激励他们。要注意他从未提及这些阴谋者的名字，以免他们变得更加无耻。他的意思如下：不是因为\Quote{你们必不会这样想}，那迷惑你们的人就可免除惩罚。他们必受惩罚；因为一人的善行不应成为另一人恶性的鼓励。说这话是为了使他们不再对他人作第二次尝试。他不仅说\Quote{那扰乱你们的人}，而且加重语气说\Quote{无论他是谁}。\par

第11节：\BibleQuote{但是，弟兄们，如果我仍宣讲割损，为什么我还受迫害？}\BibleRef{迦 5:11}\par

注意他多么清楚地将自己从指控中解脱出来，即他四处犹太化并虚伪传道。他召他们为证；因为你们知道，他说，我命令放弃法律成为他们迫害我的借口。如果我仍宣讲割损，为什么我还受迫害？这是犹太后裔能向我提出的唯一控告。如果我允许他们接受信德，同时仍保留他们祖先的习俗，那么无论信徒还是非信徒都不会陷害我，因为他们的惯常做法没有受到干扰。那么怎样！难道他没有宣讲割损吗？他难道没有给\Person{弟茂德}行割损吗？他确实行了。那么他怎能说：\Quote{我不宣讲它}？这里注意他的精确；他不是说：\Quote{我不施行割损}，而是说：\Quote{我不宣讲它}，也就是说，我不命令人如此相信。因此，不要认为这是你们教理的任何确认，因为我虽然行了割损，却没有宣讲割损。\par

第11节：\BibleQuote{那么，\OriginalTerm{十字架的绊脚石}{stumbling block of the cross}就被废除了。}\BibleRef{迦 5:11}\par

就是说，如果你们所说的属实，那障碍、那拦阻就移除了；因为连十字架对犹太人的冒犯，也不及他们祖宗的传统不应遵守的道理那么大。他们把\Person{斯德望}带到公议会前时，他们不是说这人朝拜被钉者，而是说他\Quote{攻击这圣地和法律。}\BibleRef{宗 6:13}他们也以此控告耶稣，说他违背法律。所以\Person{保禄}说，如果割损被允许，你们所陷入的纷争就平息了；此后对十字架和我们的宣讲不再有敌意。但他们为什么日复一日等待要杀害我们，却以此指控我们？是因为我把一个未受割损的人带进了圣殿\BibleRef{宗 21:29}他们才向我下手。那么，他说，我岂不是如此愚昧，在放弃了割损一事之后，还白白地、无谓地使自己遭受这些伤害，并在十字架前放下这样的绊脚石？因为你们看到，他们攻击我们，再没有什么比割损更激烈的了。那么我岂不是如此愚昧，为毫无价值的事受苦，并冒犯别人？他称之为十字架的绊脚石，因为这是由十字架的教义所命令的；而正是这一点大大冒犯了犹太人，阻碍他们接受十字架，即命令放弃他们祖宗的习俗。\par

第12节。\BibleQuote{我恨不得那搅扰你们的人，甚至将自己割绝了。}\BibleRef{迦 5:12}\par

请看，他在这里何等严厉地斥责那些欺骗者。起初，他的指责是针对那些被欺骗的人，并称他们为糊涂人，一再如此。现在，在充分纠正和训导他们之后，他转向那些欺骗者。你们也应留意他的智慧，他以前者为自己儿女，且能接受纠正的方式劝诫和惩罚他们；但对于那些欺骗者，他则视作外人且不可救药地败坏而加以断绝。他这样做，一部分是当他说：\Quote{那搅扰你们的，不论是谁，必承受惩罚。}一部分是当他发出咒诅说：\Quote{我恨不得那搅扰你们的人，甚至将自己割绝了。}他正当地说\Quote{搅扰你们}，因为那些人曾迫使他们抛弃自己的祖国、自由和天上的亲属，去寻求一个异乡和外邦的亲属；曾将他们从高处且自由的耶路撒冷逐出，迫使他们像俘虏和流放者一样漂泊。因此他咒诅他们。他的意思是：我对他们毫不关心，\BibleQuote{异端人，在第一次和第二次劝戒以后，就该弃绝他。}\BibleRef{铎 3:10}如果他们愿意，让他们不仅受割损，而且被阉割。那么，那些敢于自阉的人在哪里？他们招致宗徒的咒诅，指责天主的创造，并与摩尼派为伍。因为摩尼派称身体是背信之物，来自恶的本原；而前者以他们的行为支持这些可怜的道理，割去肢体，视之为敌意和背信。他们岂不更应该挖掉眼睛，因为欲望是通过眼睛进入灵魂的？但事实上，既不是眼睛，也不是我们身体的任何其他部分应受责备，唯有败坏的意志。但如果你不承认这一点，你为什么不因亵渎而割去舌头，因掠夺而割去双手，因恶行而割去双脚，简言之，割去全身？因为耳朵被笛声陶醉，常常使灵魂软弱；鼻孔闻到甜香，也迷惑心智，使之疯狂追求享乐。然而这将是极端的邪恶和撒殚的疯狂。那恶灵，常以杀戮为乐，诱使他们摧毁工具，仿佛造物主犯了错误，而实际上只需纠正灵魂的狂暴情慾。也许有人说：当身体被娇养时，贪慾就会炽烈，这是怎么回事？在这里也要注意，这是灵魂的罪，因为娇养肉体不是肉体的行为，而是灵魂的行为；如果灵魂选择克制肉体，它就会拥有对肉体的绝对权力。但你们所做的，就如同一个人看见有人点火、添柴、烧房子，却责怪火，而不是点火的，因为火引燃了这堆燃料并升得很高。但罪责不在火，而在点火的人；因为火被赐予，是为了煮食物、提供光明和其他类似的服务，而不是为了烧房子。同样，慾望被植入，是为了生育子女和延续生命，不是为了奸淫、通奸或淫乱；使人成为父亲，而非奸夫；成为合法的丈夫，而非诱惑者；留下继承人，而非损害他人妻子。因为奸淫不是出于本性，而是出于违反本性的放荡，本性规定使用而非滥用。我说这些不是随意的，而是作为争论的序曲，作为对那些断言天主的创造是邪恶的人的攻击；他们忽视灵魂的懈怠，疯狂地指责身体，诽谤我们的肉身；而保禄后来论及此事时，并非指责肉身，而是魔鬼般的思念。\par

第13节。\BibleQuote{弟兄们，你们蒙召是为了自由；但不要将这自由作为放纵肉慾的机会。}\BibleRef{迦 5:13}\par

从此，他似乎转入道德论述，却以新的方式，在他的其他书信中从未出现。因为所有书信都分为两部分：第一部分讨论教义，最后部分讨论生活准则；但在此处，他进入道德论述后，又将其与教义部分结合起来。因为这段经文在反对摩尼教徒的争论中涉及教义。\Quote{你们不可将这自由当作放纵肉欲的机会？}是什么意思？他说，基督已将你们从奴役的轭下解救出来，祂让我们自由行事，并非要我们滥用自由去作恶，而是要我们获得更高赏赐的根据，进到更高超的哲学。因他屡次称法律为奴役之轭和诅咒的缘由，为避免有人怀疑他倡导废弃法律是让人过无法律的生活，他纠正这观念，声明其目的不是让我们的生活无法度，而是让我们的哲学超越法律。因为法律的束缚已被打破，我这样说，不是为降低标准，而是为提升它。因为行淫者和守贞者都越过了法律的界限，但方向不同：前者被引向更坏，后者被提升到更好；前者违反法律，后者超越法律。因此保禄说，基督从你们身上除去了轭，不是要你们蹦跳踢蹬，而是要你们虽无轭却以均匀的步伐前行。接着他说明容易实现这目标的方法。是什么方法呢？他说：\par

第 13 节。\Quote{反倒要藉着爱德彼此服事。}\newline{}\newline{}\BibleRef{迦 5:13}\par

此处他又暗示，争端、派系精神、统治欲和狂妄自大是他们错误的根源，因为统治欲是异端之母。他说\Quote{彼此服事}，表明这邪恶源于这种狂妄傲慢的精神，因此他应用相应的补救。正如你们的纷争来自彼此主宰的欲望，\Quote{你们要彼此服事}；这样你们将再次和好。然而，他没有公开指责他们的过失，却公开告诉他们纠正之道，借此他们能意识到那过失；犹如有人不直接指出放荡者的放荡，却不断劝他向洁德努力。爱邻人如己者，不会拒绝比任何仆人更谦卑地服事他。如同火接触蜡，轻易将其软化，爱的温暖比火更有力地消融一切傲慢和狂妄。因此他不只说\Quote{彼此相爱}，而是说\Quote{彼此服事}，以此表明情感的热烈。当他们被解除法律的轭，免于脱缰跑开时，另一个轭被加上——爱的轭，比前者更强，却更轻省愉快；为指出服从它的方法，他加上：\par

第 14 节。\Quote{因为全部法律总括在这句话内：\InnerQuote{你应当爱你的近人如你自己。}}\newline{}\newline{}\BibleRef{迦 5:14}\par

看到他们如此看重法律，他说：\Quote{你们若愿意成全法律，就不要受割损，因为法律不是成全在割损上，而是成全在爱上。}请注意，他无法忘记他的忧伤，即使进入道德论述，也经常触动困扰他的事。\par

第 15 节。\Quote{但如果你们彼此咬噬、并吞，你们该小心，免得彼此消灭。}\newline{}\newline{}\BibleRef{迦 5:15}\par

为不使他们伤心，他虽明知实情，却不直说，而是含糊地提及。因为他不说\Quote{既然你们彼此咬噬}，也未在接下来的句子中断言他们将彼此消灭，而是说\Quote{你们该小心，免得彼此消灭}，这是忧虑和警告的语气，而非定罪。他所用的词语意味深长：他不仅说\Quote{你们咬噬}（这在冲动时可能发生），而且说\Quote{并吞}，这含有怀恨之意。咬噬是发泄怒气，但并吞是极度凶残的证明。他所说的咬噬和并吞并非肉体上的，而是更加残忍的；品尝人肉所造成的伤害，远不及把牙嵌入别人灵魂的伤害。灵魂比肉体贵重多少，对其的损害就更严重。\Quote{免得彼此消灭}。因为行恶和设谋的人，为的是消灭别人；因此他说：小心这邪恶不要落在你们自己头上。因为斗争和纷争，对接纳它们和引进它们的人，都是毁灭和破坏，比蛀虫更彻底地吞噬一切。\par

第 16 节。\Quote{我告诉你们：你们若随从圣神引导，就绝不能去满足肉身的纵欲。}\newline{}\newline{}\BibleRef{迦 5:16}\par

此处他指另一条使责任变得容易、并确保所说之事的道路，一条产生爱、且被爱围护的道路。因为没有什么比属神更能使我们易于去爱，也没有什么比爱的力量更能促使圣神住在我们内。因此他说：\Quote{你们若随从圣神引导，就绝不能去满足肉身的纵欲。}他说了疾病的成因后，也提及带来健康的药方。这药方是什么？我们所说诸恶的毁灭是什么？就是随从圣神的生活。因此他说：\Quote{你们若随从圣神引导，就绝不能去满足肉身的纵欲。}\par

第 17 节。\Quote{因为肉身的纵欲相反圣神，圣神的引导相反肉身。二者互相敌对，致使你们不能行你们所愿意的事。}\newline{}\newline{}\BibleRef{迦 5:17}\par

这里有人指责宗徒把人分成了两部分，并称他所构成的本性与自身相冲突，身体与灵魂争斗。但这绝然不是如此；因为他所谓\Quote{肉性}，并非指身体；若是指身体，紧接着的句子\Quote{因为它渴望}，他说，\Quote{反对圣神}是什么意思？然而身体不动，是被动的，不是主动者，而是受动者。那么它如何渴望呢？因为渴望属于灵魂，不属于身体，因为在另一处经上说：\BibleQuote{我的灵魂渴慕}\BibleRef{咏 83:2}，又\BibleQuote{凡你灵魂所愿望的，我必为你成就}\BibleRef{撒上 20:4}，又\Quote{不要随你心里的欲望行走}，\BibleQuote{我的灵魂渴慕}\BibleRef{咏 41:1}。那么为什么保禄说\Quote{肉性渴望反对圣神}呢？他惯常称肉性，不是指本性身体，而是指堕落的意志，正如他说：\BibleQuote{但你们已不在肉性内，而在圣神内}\BibleRef{罗 8:8}，又说：\BibleQuote{凡属肉性的人，不能中悦天主}。那么怎样？肉性要被消灭吗？那这样说的人不是也穿着肉身吗？这样的道理不是出于肉性，而是出于魔鬼，因为\BibleQuote{他从起初就是杀人者}\BibleRef{若 8:44}。那么他的意思是什么？他所称的肉性是指世俗的心思、懒散和疏忽，这不是对身体的指控，而是对懒散灵魂的控诉。身体是一个工具，没有人会对工具感到厌恶和憎恨，而是对滥用它的人。因为我们憎恨和惩罚的不是铁器，而是杀人者。但也许有人说，把灵魂的过错称为肉性本身就已经是对身体的指控。我承认肉性低于灵魂，但它也是好的，因为低于好的东西本身可以是好的，但恶不是低于好，而是与它对立的。如果你能向我证明恶起源于身体，你就可以随意指控它；但如果你试图把它的名字转变为对它的指控，你也应该同样指控灵魂。因为失去真理的人被称为\BibleQuote{属血气的人}\BibleRef{格前 2:14}，而魔鬼的族类被称为\BibleQuote{邪恶的鬼神}\BibleRef{弗 6:12}。\par

再者，圣经惯常将\Quote{肉性}之名给予圣体奥迹，并给予整个教会，称它们为基督的身体。\BibleRef{哥 1:24} 不仅于此，为了使你视那些以肉体为媒介的事物为祝福，你只需想象感官的熄灭，就会发现灵魂失去了一切辨识力，对其先前所知一无所知。因为如果天主的德能\BibleQuote{自从创世以来，借着所造之物就被看得清清楚楚}\BibleRef{罗 1:20}，没有眼睛我们怎能看见它们？如果信德\BibleQuote{来自听闻}\BibleRef{罗 10:17}，没有耳朵我们怎能听见？宣讲取决于巡回，其中用到舌头和脚。因为\BibleQuote{若不被派遣，怎能宣讲呢？}\BibleRef{罗 10:15} 同样，写作是靠手进行的。你没有看到肉体的服务为我们带来千百种益处吗？在他所说\Quote{肉性渴望反对圣神}的表述中，他是指两种精神状态。因为这些是彼此对立的，即德行与恶行，而非灵魂与身体。若后两者如此对立，它们就会彼此毁灭，如同火与水，光明与黑暗。但如果灵魂照顾身体，为其深谋远虑，忍受千百种事情以免离开它，并抗拒与之分离，而身体也为灵魂服务，向其传达许多知识，并适应其运作，它们怎能是相反且彼此冲突的呢？就我而言，我从它们的行为看出，它们不但不相反，反而密切协调、彼此依附。因此，他所说的对立不是指这些，而是指善恶原则的争战。\BibleRef{罗 7:23} 意愿与不意愿属于灵魂；因此他说：\Quote{这些是彼此相反的}，使你不要任凭灵魂在其邪恶欲望中前行。因为他说话如同一个主人和老师，带着威吓的方式。\par

第十八节。\BibleQuote{但如果你们被圣神引导，你们就不在法律之下。}\BibleRef{迦 5:18}\par

如果问这两者如何相连，我回答：密切而明显；因为拥有圣神的人，若恰当拥有，就借此熄灭一切邪恶欲望，而摆脱这些的人不需要法律的帮助，反而远超其诫命。从不发怒的人，何需听\Quote{不可杀人}的命令？从不放荡眼神的人，何需训诫\Quote{不可奸淫}？谁会与拔除了恶行之根的人谈论恶果？因为愤怒是杀人的根，而窥视面孔是奸淫的根。因此他说：\Quote{如果你们被圣神引导，你们就不在法律之下。}这里他好像发表了对法律的一篇崇高而显著的颂辞，如果至少法律按其能力，在圣神降临前代替了圣神的位置。但我们不必因此继续与我们的启蒙老师一起。那时我们公正地服从法律，为了用恐惧克制我们的欲望，圣神尚未显现；但现在恩宠已赐下，它不仅命令我们戒除这些，而且还熄灭它们，并引导我们走向更高的人生规则，那么法律还有什么需要呢？从内在冲动达到崇高卓越的人，不需要启蒙老师，就像没有哲学家需要语法老师。那么你们为何如此贬低自己，以前已给了圣神，现在却仍听从法律？\par

\Quote{第19至21节：}\BibleQuote{邪淫、污秽、淫荡、崇拜偶像、行邪术、仇恨、争执、嫉妒、忿怒、分党、分裂、异端、嫉妒、酗酒、狂欢，以及类似的事。关于这些，我从前警告过你们，如今再警告你们：行这样事的人，不得继承天主的国。}\BibleRef{迦 5:19-21}\par

现在回答我，你这控告自己肉身的人，你认为这话是对肉身说的，如同对待敌人和对手。假定奸淫和邪淫如你所说，是从肉身而来；然而仇恨、争端、嫉妒、纷争、异端和巫术，这些仅仅来自败坏的道德选择。其他也是如此，因为它们怎能属于肉身呢？你们发现，他在这里不是在谈论肉身，而是谈论地上的思想，它们紧贴地面。因此他也用\BibleQuote{行这样事的人，不得继承天主的国}来警告他们。如果这些属于本性，而不是败坏的道德选择，那么他的说法\Quote{行}就不恰当，应该说\Quote{受}。为什么他们要被逐出天国？因为赏罚不是针对来自本性的，而是来自选择的。\par

\Quote{第22节：}\BibleQuote{然而圣神的效果是仁爱、喜乐、平安。}\BibleRef{迦 5:22}\par

他说不是\Quote{圣神的作为}，而是\Quote{圣神的效果}。然而，灵魂是多余的吗？提到了肉身和圣神，但灵魂在哪里？他是在谈论没有灵魂的存在吗？因为如果肉身的事是邪恶的，圣神的事是良善的，那么灵魂必定是多余的。绝非如此，因为制服私欲偏情属于她，并与她有关；她处于恶习和美德之间，如果她恰当地使用了身体，便使身体成为属神的；但如果她与圣神分离，沉溺于邪恶的欲望，她便使自己更加属世。你们始终发现，他的论述不是关于肉身的本质，而是关于道德选择，即其是否有恶习。为什么他说\Quote{圣神的效果}呢？因为恶行只源于我们自己，因此他称它们为\Quote{作为}；而善行不仅需要我们的努力，也需要天主的慈爱。他先列出这些善事的根源，然后逐一数算它们，用这些话：\BibleQuote{仁爱、喜乐、平安、忍耐、良善、温和、忠信、柔和、节制；这样的事没有法律禁止。}因为谁会向一个自身拥有一切，并以爱为完美哲学的主宰的人发号施令呢？正如驯良的马，自动做一切事，不需要鞭打；同样，那藉着圣神达致卓越的灵魂，也不需要法律的训诫。在此他也彻底而有力地摒弃了法律，不是因为它不好，而是因为它低于圣神所赐的哲学。\par

\Quote{第24节：}\BibleQuote{凡属于基督耶稣的人，已把肉身连同私欲偏情钉在十字架上了。}\BibleRef{迦 5:24}\par

为了避免他们反驳：\Quote{谁是这样的人呢？}他通过他们的行为指出那些达到这完美境界的人，在此再次将\Quote{肉身}之名用于恶行。他并不是说他们毁灭了肉身，否则他们怎么活呢？因为被钉在十字架上的东西是死的、不起作用的，但他指示了完美的生活准则。因为私欲，尽管烦扰，却徒然发怒。既然圣神有这样的能力，我们就该活在圣神内，并以此为满足，正如他本人补充说：\par

\Quote{第25节：}\BibleQuote{如果我们因圣神生活，就也当随从圣神而行。}\BibleRef{迦 5:25}\par

——受祂的法则所管理。因为这是\Quote{行走}一词的力量，也就是说，我们当以圣神的能力为满足，不再从法律寻求帮助。然后，指出那些想要引入割礼的人是被野心所驱动，他说：\par

\Quote{第26节：}\BibleQuote{我们不要贪图虚荣，}这是一切邪恶的起因，\BibleQuote{彼此惹气，}以致争执和纷争，\BibleQuote{彼此嫉妒，}因为虚荣产生嫉妒，嫉妒产生这无数邪恶。\BibleRef{迦 5:26}\par

\SourceRecord{Translated by Gross Alexander.；Nicene and Post-Nicene Fathers, First Series；Vol. 13.；Edited by Philip Schaff.；Buffalo, NY: Christian Literature Publishing Co.,；1889.；<http://www.newadvent.org/fathers/23105.htm>.}

% structure: p0001 source=h1 target=section reason=page-title

\section{第六篇讲道稿 论迦拉达书}

第一节 \BibleRef{迦 6:1}\par

\BibleQuote{\Emph{弟兄们}, \Emph{即使有人被任何过犯所胜。}}\par

\Emph{鉴于}在谴责的掩饰下，他们满足了自己的私情，并以纠正已犯的过失为名，实则推进自己的野心，保禄说：\BibleQuote{弟兄们，若有人被\OriginalTerm{所胜}{overtaken}}。他没说\Quote{若有人犯}，而是说若被\OriginalTerm{所胜}{overtaken}，即被卷入。\par

\BibleQuote{你们属神的人，要把这样的人挽回过来}\par

他没说\Quote{惩罚}，也没说\Quote{审判}，而是说\Quote{纠正}。他并不止于此，为了表明他们应当以极温和的态度对待那些失足的人，他补充道：\par

\BibleQuote{以温柔的心}\par

他没说\Quote{以温柔}，而是说\Quote{以温柔的心}，以此表明这是蒙圣神悦纳的，并且能以温和来施行纠正乃是一种神恩。然后，为了防止那纠正别人的人因自己的职责而过分自高，他用同样的畏惧来提醒他，说：\par

\BibleQuote{也当小心，免得你也受诱惑}\par

因为正如富人向穷人输送捐助，唯恐自己一旦陷入贫穷也能得到同样的周济，我们同样也应该这样做。因此，他提出这个有力的理由，说：\Quote{也当小心，免得你也受诱惑}。他首先以\Quote{若有人被\OriginalTerm{所胜}{overtaken}}为犯罪者辩护；其次，用一个表示极大软弱的词；最后，以\Quote{免得你也受诱惑}这句话，将罪责归于魔鬼的恶意，而非灵魂的懈怠。\par

第二节 \BibleRef{迦 6:2} \BibleQuote{你们要彼此担待重担}\par

既然人不可能没有过犯，他就劝告他们不要严苛地审视他人的过错，而是要甚至担当他们的软弱，这样他们自己的软弱也可以由他人担当。正如建造房屋时，并非所有石头都处于同一位置：有的适合做房角石却不适合做地基，有的适合做地基却不适合做房角石。在教会身体上也是如此。我们肉体的结构也是一样；然而，一个肢体担当另一个肢体，我们并不向每个肢体要求一切，而是每个肢体所共同贡献的构成了身体和建筑。\par

第二节 \BibleRef{迦 6:2} \BibleQuote{这样就成全了基督的律法}\par

他没说\Quote{成全}，而是说\Quote{完成}；意思是，你们众人共同来补足它，通过彼此担待的事。例如：这人易怒，你性子慢；那么你担当他的急躁，好让他也担当你的迟钝；这样，既然你支持他，他不会越界；既然弟兄容忍你，你在自己软弱之处也不会犯罪。所以，你们彼此伸手扶持将要跌倒的人，就共同成全了律法，各人以自己的忍耐补足邻人的不足。但若不这样，而是每人追究邻人的过失，那么你们将一事无成。如同身体，若要求每个肢体发挥同一功能，身体就无法存在；同样，如果我们向众人要求一切，弟兄之间必有大纷争。\par

第三节 \BibleRef{迦 6:3} \BibleQuote{人若以为自己是什么，其实什么也不是，就是欺骗自己}\par

这里他又斥责他们的骄傲。那以为自己是什么的人，其实什么都不是，并且从一开始就以此心态证明了自己的一无是处。\par

第四节 \BibleRef{迦 6:4} \BibleQuote{但各人要察验自己的行为}\par

这里他表明我们应当省察自己的生命，并且不是轻率地，而是仔细地衡量自己的行为；例如，若你行了一件善事，要思考是否出于虚荣，或出于必要，或出于恶意，或出于虚伪，或出于其他属人的动机。因为正如金子在被放入熔炉之前显得明亮，但投入火中后就被彻底检验，一切掺假的都与纯金分离；同样，我们的行为若经仔细查验，就会显明，我们会发现自己受到了许多责备。\par

第四节 \BibleRef{迦 6:4} \BibleQuote{这样，他就只有在自己身上夸口，而不在别人身上}\par

他说这话，不是作为规定，而是作为让步；他的意思是：夸耀是愚昧的，但如果你要夸耀，不要像法利塞人那样夸耀自己的邻人。因为这样受教的人很快就会完全放弃夸耀；因此他让步一部分，以便逐渐铲除全部。那只在自己身上夸口、不拿别人相比的人，很快就会改正这个缺点。因为他不认为自己比别人好——这就是\Quote{不在别人身上}的意思——而是只通过省察自己而自高，之后便会停止这样做。为了使你确信这就是他想要确立的，请注意他如何以恐惧来抑制他，前面说\Quote{各人要察验自己的行为}，并在此补充道：\par

第五节 \BibleRef{迦 6:5} \BibleQuote{因为各人要担当自己的担子}\par

他似乎在提出一个禁止夸耀别人的理由；但同时他纠正了夸耀者，使他不再自高，通过让他想起自己的过失，并把\OriginalTerm{担子}{burden}和重负的意念压在他的良心上。\par

第六节 \BibleRef{迦 6:6} \BibleQuote{但那位在圣言上受教的，应让施教的人分享一切美物}\par

他接着论及教师，指出门徒应殷勤照料他们。那么基督为何如此命令呢？因为这条法律，\BibleQuote{传福音的靠福音养生}，\BibleRef{格前 9:14}是在新约中规定的；同样在旧约中，\BibleRef{户 31:47}肋未人也从百姓那里获得许多进项。那么，我问，祂为何如此安排？岂不是为了预先奠定谦卑和爱的基础吗？因为教师的尊位往往使拥有它的人骄傲，为了抑制他的精神，祂便强加给他一个必要，即从门徒手中请求帮助。反过来，祂又赐给门徒培养善意的方法，借着要求他们对教师所行的善举，来训练他们对他人也要温柔。这样，双方就产生了不少的感情。若不是我所陈述的原因，那曾用玛纳喂养迟钝的犹太人的祂，为何使宗徒们落到请求帮助的必要中呢？这岂不明显地表明，祂的目标是谦卑和爱的巨大益处，并要使受教者不以表面卑微的教师为耻吗？请求帮助带有耻辱的外表，但当他们的教师大胆地提出要求时，这耻辱便消失了，门徒们由此获得了不少益处，学会了藐视一切外表。因此他说：\BibleQuote{但那在圣言上受教的，应让施教者共享一切美物。} 这就是说，应对他显示全部慷慨；他用\InnerQuote{一切美物}这个词来表达。他说，门徒不应保留任何为自己，而应一切公用，因为他所收受的比他所给予的更好——正如天上之物优于地上之物。他在另一处表达了这一点：\BibleQuote{我们若给你们播下了属神的事，若收取你们属血肉的事，还算大事吗？}\BibleRef{格前 9:11}因此，他称这一程序为\Quote{通传}，表明一种交换发生了。爱也借此大大地受到促进和坚固。如果教师只要求适当的限度，他接受帮助并不会损害自己的尊严。因为这是值得称赞的，如此勤勉地致力于圣言，以致需要他人的帮助，处于多种贫穷之中，并漠视一切生计。但如果他超越适当限度，他损害了尊严，不是仅仅因为接受，而是因为接受过多。那么，为了避免教师的恶行使门徒在这件事上更加懈怠，并因他的行为而常常忽略他（尽管他很贫穷），他接着说：\par

第9节。\BibleQuote{我们行善不要厌倦。}\BibleRef{迦 6:9}\par

这里，他指出了这种野心与世俗事务之间的区别，说：\BibleQuote{不要错误；天主不是好嘲弄的；因为人种什么，就收什么。那顺着自己肉身播种的，必从肉身收获败坏；但那顺着圣神播种的，必从圣神收获永生。} 如同种子的情况，种豆子的不能收谷子，因为所种的和所收的必须同类；在行动上也是如此：在肉身上种植纵情、醉酒或不节制的欲望的人，必收获这些事的果实。这些果实是什么？惩罚、报应、耻辱、嗤笑、毁灭。因为丰盛的餐桌和菜肴，结局无非是毁灭；它们自己败坏，也毁灭身体。但圣神的果实，其性质与这些不同，而且全然相反。试想：你播种了施舍？天上的宝库和永恒的光荣等着你；你播种了节制？尊荣、赏报、天使的掌声，以及从审判者而来的冠冕等着你。\par

第9、10节。\BibleQuote{我们行善不要厌倦；因为若不灰心，到了适当的时候，我们就要收获。所以，我们一有机会，就应向众人行善，尤其应向有同样信仰的家人。}\par

为避免有人以为只有他们的老师才应受照顾和供养，而其他人则可以忽略，他便将话题普遍化，为此种慈善的热忱向众人敞开大门；不仅如此，他甚至将此推到如此高度，以致命令我们向犹太人和希腊人都应显示慈爱，固然是按照适当的次序，但仍要显示慈爱。这次序是什么？就是对信友给予更大的关心。他在这里的努力与在其他书信中相同；他不仅谈论显示慈爱，且要怀着热忱和恒心去做，因为\Quote{播种}和\Quote{不倦怠}的说法正暗示了这一点。然后，在要求一项伟大工作之后，他将其赏报放在眼前，并提到一种新奇而奇妙的收获。在农夫中，不仅是播种者，连收割者也忍受许多劳苦，要与干旱、尘土和艰辛劳作搏斗，但在此情况下，这些都不存在，正如他借着\Quote{因为到了适当的季节，我们必会收获，如果我们不倦怠}这句话所指明的。他用这种方式激励并吸引他们；他又用另一种动机催促和推动他们，说：\Quote{所以，我们一有机会，就应行善。}正如播种不总是在我们的能力之内，显示慈爱也是如此；因为我们一旦被带离这里，即使我们千百次愿意，也将再不能成就什么。为我们的这个论点，十童女\BibleRef{玛 25:1}可以作证，她们虽然千百次愿意，却被关在新郎外面，因为她们没有带来丰富的爱德。那忽略了拉匝禄的富翁\BibleRef{路 16:19}也是如此，因为他缺乏这种援助，虽然哭泣并多次恳求，却没有从圣祖或任何人那里得到怜悯，反而一直得不到任何宽恕，并受永火折磨。因此他说：\Quote{我们一有机会，就当向众人行善，}借此尤其将他们从犹太人的狭隘心思中解放出来。因为他们的全部仁爱只限于本族，但恩宠所赐的生活准则却邀请陆地和海洋共赴慈善的筵席，只是对自己的家人表现出更大的关心。\par

第11-12节。\BibleQuote{你们看，我亲手写给你们的字是何等的大。那些希图在肉体上炫耀的人，他们强迫你们受割损。}\par

请注意这蒙福的灵魂为怎样的忧伤所占据。如同那些被某种悲伤压倒的人，他们失去了自己的亲人，遭受了意想不到的灾祸，日夜不得安息，因为忧伤围困着他们的灵魂，同样，蒙福的保禄在短暂的道德训诲之后，又回到那先前主要扰乱他心思的主题上，说如下的话：\Quote{你们看，我亲手写给你们的字是何等的大。}借此他表明整封书信是他亲自写的，这是极大真诚的证明。在他的其他书信中，他只是口述，由别人书写，正如从\WorkTitle{罗马书}明显可见，因为在它的结尾处说：\BibleQuote{我德尔图斯，这封书信的执笔者，问候你们；}\BibleRef{罗 16:22}但在这一次，他全部亲自书写。他这样做是出于必要，而不仅出于情感，乃是为了消除一个有害的猜疑。他被指控参与了一些与他无关的行为，并被传言说他在宣讲割损，却假装不宣讲，因此他被迫亲手写下这封书信，从而预先留下书面的证据。按我的理解，通过\Quote{何等的大}这个说法，他不是指字母的大小，而是指形状的丑陋，好像他在说：\Quote{我虽然不善于书写，却被逼亲手写信，以堵住这些诽谤者的口。}\par

第12-13节。\BibleQuote{那些希图在肉体上炫耀的人，他们强迫你们受割损，只是为避免因基督的十字架而受迫害。因为连那些受割损的人自己也不遵守法律；但他们却愿意你们受割损，好为你们的肉体而夸耀。}\par

这里他指出，他们如此忍受，不是出于自愿，而是出于必要，并为他们提供了一个退路的机会，几乎是在为他们辩护，劝他们尽快离开自己的老师。\Quote{在肉体上炫耀？}是什么意思？意思是被人尊重。因为他们因离弃祖宗的习俗而受犹太人辱骂，他说，他们想要伤害你们，使他们自己不至于被指控这一点，而是借着你们的肉体为自己辩护。他此处的目的是要表明，他们如此行不是出于对天主的敬畏；仿佛他在说，这种做法不是基于虔敬，这一切都是出于人的野心；为使不信者因信友的伤残而满意，他们宁愿得罪天主以取悦人；这就是\Quote{在肉体上炫耀}的意思。然后，作为证明他们因另一个原因也是不可宽恕的，他再次使他们相信，他们命令此事不仅为了取悦他人，也为了自己的虚荣。因此他补充说：\Quote{好为你们的肉体而夸耀，}好像他们是弟子和老师似的。而这是什么证明呢？他说：\Quote{因为连他们自己也不遵守法律；}即使他们遵守，也会招致严厉的谴责，但现在他们的目的本身就是败坏的。\par

第14节。\BibleQuote{至于我，我决不夸耀，只夸耀我们的主耶稣基督的十字架。}\BibleRef{迦 6:14}\par

这个标记确实被视为可鄙的；但这只是在世人眼中如此，在天上及信徒中却是至高荣耀。贫穷也可鄙，却是我们的夸耀；被人轻看使他们发笑，我们却因此振奋。同样，十字架是我们的夸耀。他没有说：\Quote{我不夸耀，}也没有说：\Quote{我不会夸耀，}而是说：\Quote{我绝不夸耀，}好像他厌恶它如同荒谬之事，并祈求天主的帮助以在其中成功。十字架的夸耀是什么？就是基督为了我取了奴仆的形象，为我这奴仆、仇敌、无情的人忍受苦难；是的，祂如此爱我，以至为我舍弃自己，甚至受诅咒。还有什么可相比的！如果奴仆仅因得主人称赞而兴奋（他们与主人在本性上相近），那么当主人——真实的天主——不以为我们所受的十字架为耻时，我们怎能不夸耀呢？让我们不要以祂无可言喻的柔情为耻；祂不以为你被钉十字架为耻，而你却羞于承认祂无限的关怀？这就像一个囚犯，不曾以他的国王为耻，但当国王来到监狱亲自为他解开锁链时，他反而因此羞耻。但这简直是极致的疯狂，因为正是这一事实成为夸耀的特殊理由。\par

14节：\BibleQuote{藉此，世界已被钉死于我，我亦被钉死于世界。}\BibleRef{迦 6:14}\par

此处他所说的世界并非天或地，而是生活中的事务、人的赞誉、随从、荣耀、财富，以及一切看似辉煌之物。这些都于我已成死物。基督徒应当如此，并常说这样的话。他不仅满足于前述的死，还加上另一句：\Quote{我也于世界，}意指双重的死，并说：它们于我已是死的，我于它们也是死的；它们再也不能俘虏和胜过我，因为它们已一次而死；我也不能再贪恋它们，因为我也于它们死了。没有比这死更有福的，因为它是有福生命的基础。\par

15、16节：\BibleQuote{原来割损算不得什么，不割损也算不得什么，而是一个新受造物。凡以此准则而行的，愿平安与怜悯临于他们，也临于天主的以色列。}\par

请观察十字架的能力，它将他提升到何等高度！不仅为他将一切世俗之事置于死地，更使他远超旧约。什么能与这能力相比？十字架已说服他（他原愿为割损而杀人并被杀）将割损与不割损等量齐观，并寻求新奇、奇妙、超越诸天之事。他把我们的生活准则称为\Quote{新受造物，}既因过去也因将来；因过去，我们的灵魂曾被罪的旧态所衰老，却藉着圣洗立时更新，如同再造。因此我们需要一种新的、天上的生活准则。至于将来，因为天、地及一切受造物，都将随我们的身体转化为不朽。因此，不要向我提割损，它现已毫无用处（当万物都经历如此变化时，它又如何显现？）而要寻求恩宠的新事。因为追求这些事的人将享有平安与和睦，并可恰当地被称为\Quote{以色列}之名。而持相反意见的人，即使是以色列的后裔并拥有其名号，也完全失落了这一切，包括关系和名号本身。但那些遵守此准则、远离旧路、追随恩宠的人，才有能力成为真正的以色列人。\par

17节：\BibleQuote{从今以后，不要再有人麻烦我了。}\BibleRef{迦 6:17}\par

他说这话并非因疲倦或被压倒；他本愿为门徒的缘故行事并受苦；他曾说：\BibleQuote{务要专心，不论时机是否合宜，}\BibleRef{弟后 4:2}又说：\BibleQuote{或许天主会赐给他们悔改，得以认识真理，并摆脱魔鬼的网罗，}\BibleRef{弟后 2:25}他怎么会现在松懈倒退呢？他说这话究竟为何？是为了激励他们懈怠的心思，使他们更深刻地敬畏，并坚定他所制定的法规，抑制他们不断的动摇。\par

17节：\BibleQuote{因为我身上带着耶稣的烙印。}\BibleRef{迦 6:17}\par

他没有说\Quote{我有，}而是说\Quote{我带着，}如同一个人以战利品和王旗为荣。虽然细想之下看似耻辱，但这人却夸耀自己的伤口，犹如军旗手一般，他得意地带着这些伤口。他为何这样说？\Quote{比起任何论证、任何言语，这些伤口更能为我辩护，}他说，\Quote{因为这些伤口所发出的声音比号角更响亮，反对我的人，以及那些说我是在教导中演戏、讨人欢心的人，都在此无声。因为看见一个士兵满身血迹、带着千疮百孔从战场退下的人，绝不敢指责他懦弱和背叛，因为他身上带着他勇气的明证。所以你们也该这样判断我。如果有人想听我的辩护并了解我的观点，就让他看看我的伤口吧，这些比言辞和书信更强有力的证明。在他的书信开头，他以突然的皈依显明了真诚；在结尾，他以皈依所伴随的危险来证明。为避免有人指责他并非故意改变方向（虽然意向正直），但没有坚持同一目的，他便以他的试炼、危险和鞭痕作见证，证明他确实坚持了。}\par

然后，他在每一细节上清楚证明了自己清白，并证实他所说的一切并非出于愤怒或恶意，而是对他们的感情毫无损伤，于是他以一个满含千种祝福的祈祷来结束讲辞，以此再次确认这同一观点，祈祷词如下：\par

第18节：\BibleQuote{愿我们的主耶稣基督的恩宠与你们的心灵同在，弟兄们。阿们。}\BibleRef{迦 6:18}\par

藉此最后一句，他封印了前面所说的一切。他不只是说\Quote{与你们同在}，如别处所说，而是说\Quote{与你们的灵同在}，以此将他们从属肉体的事物中拉出，并处处彰显天主的恩惠，提醒他们所享有的恩宠，从而得以将他们从一切犹太化的错误中召回。因为领受圣神并非来自法律的贫乏，而是来自由信德而来的义德；而领受之后保守它，也并非来自割损，而是来自恩宠。为此，他以一个祈祷来结束劝言，提醒他们恩宠和圣神，同时称呼他们为弟兄，并恳求天主使他们能继续享有这些恩惠，从而为他们提供了双重的保障。因为祈祷和教导都趋向同一件事，并一同成为他们的双重围墙。因为教导提醒他们享受了哪些恩惠，从而更能使他们持守教会的道理；而祈祷呼求恩宠，劝勉他们恒心持守，不容圣神离开他们。祂住在他们内时，他们所持守的这些学说的所有错误，就都像灰尘一样被抖落了。\par

\SourceRecord{Translated by Gross Alexander.；Nicene and Post-Nicene Fathers, First Series；Vol. 13.；Edited by Philip Schaff.；Buffalo, NY: Christian Literature Publishing Co.,；1889.；<http://www.newadvent.org/fathers/23106.htm>.}
